% PDF 247--302；原书第 219--274 页。
\chapter{大语言模型的挑战与应对方法}
\label{ch:challenges}
\chaptercredit{6}
\chapterabstract{大语言模型带来了一些特有的挑战。要以合乎伦理且有效的方式部署这些模型，就必须应对这些挑战。本章关注大语言模型的幻觉、毒性、偏见、公平性和隐私问题。首先讨论幻觉，即大语言模型生成事实不正确或毫无意义的内容，并介绍衡量和减少幻觉的方法。接着讨论毒性、偏见和公平性，考察大语言模型如何在无意间延续刻板印象或不公正待遇，以及如何评估和纠正这些偏差。随后讨论隐私问题，指出大语言模型在大量数据上训练所带来的风险，并介绍保护用户隐私的技术。每节都概述相应挑战、成因、用于衡量和量化的指标、评估基准以及当时的缓解策略。本章最后用一个实践教程演示如何在大语言模型中使用偏见缓解技术，说明它们对模型行为的影响。}

% PDF 247--257；原书第 219--229 页。
\section{幻觉}
\label{sec:ch06-hallucination}

\begin{keypoint}
语言模型的幻觉（hallucination）是指生成的文本虽然语法正确、表面连贯，却偏离事实或原始材料的意图。幻觉对大语言模型的可靠性及其在实际应用中的使用构成很大挑战。
\end{keypoint}

\refcite{ch06-ref-129}{Zhang 等人（2023）}的综述将大语言模型的幻觉主要分为三类。

\begin{enumerate}
\item \textbf{与输入冲突的幻觉（input-conflicting hallucination）。}大语言模型的输出偏离用户提供的输入，其中通常包括任务指令和实际待处理的内容。这类偏差可能表明模型误解了用户意图，或未忠实处理所给内容。机器翻译和摘要等传统自然语言处理任务中也常见这类问题。
\item \textbf{与上下文冲突的幻觉（context-conflicting hallucination）。}在较长的对话或多轮交互中，大语言模型可能给出前后不一致的回答。这表明模型跟踪上下文或在长交互中保持连贯的能力可能不足，原因可能在于难以保留记忆或辨识上下文。
\item \textbf{与事实冲突的幻觉（fact-conflicting hallucination）。}大语言模型生成与已知事实相悖的内容时，就会出现这类幻觉。它的成因多样，可能出现在模型生命周期的不同阶段。例如，询问历史事实时，大语言模型可能给出错误答案，传播错误信息。由于没有一套确切的知识库可供参照，与事实冲突的幻觉是三类幻觉中最复杂的问题。这类幻觉往往对大语言模型实际部署的影响更广，因此受到研究者更多关注。
\end{enumerate}

\subsection{成因}
研究者考察大语言模型的幻觉时，发现以下几种促成因素。

\begin{enumerate}
\item \textbf{知识不足。}大语言模型可能缺少必要知识，或在预训练中吸收了错误信息。原因是它对本体知识的记忆和推理能力并不完善，见 \refcite{ch06-ref-061}{Li 等人（2022）}和 \refcite{ch06-ref-118}{Wu 等人（2023）}的研究。这些不足可能使大语言模型以过强的把握给出编造的回答。
\item \textbf{训练数据偏差。}训练数据的性质会影响大语言模型产生幻觉的倾向。\refcite{ch06-ref-074}{McKenna 等人（2023）}发现，幻觉与训练数据分布有关，尤其是在模型倾向于确认训练集所支持的假设时。
\item \textbf{人类语料可能有误。}人生成的语料也可能不准确，含有过时、带有偏差或编造的内容。在这些数据上训练的大语言模型，很可能在输出中复现这些错误（\refcite{ch06-ref-016}{Chang et al., 2019}；\refcite{ch06-ref-024}{Dziri et al., 2022}；\refcite{ch06-ref-065}{Liska et al., 2022}；\refcite{ch06-ref-090}{Penedo et al., 2023}）。
\item \textbf{回答过于自信。}大语言模型常常高估自己的知识范围，给出过于自信而且错误的回答。\refcite{ch06-ref-051}{Kadavath 等人（2022）}和 \refcite{ch06-ref-122}{Yin 等人（2023）}指出了这一问题，即使 GPT-4 等先进模型的表现，与人类表现的参照也有明显差距。
\item \textbf{对齐与迎合。}预训练后的对齐过程若没有考虑大语言模型已有的知识，可能诱发幻觉（\refcite{ch06-ref-099}{Schulman, 2023}）。这种偏差再加上大语言模型倾向于附和用户观点，即迎合（sycophancy），会使问题加剧（\refcite{ch06-ref-091}{Perez et al., 2022}）。
\item \textbf{生成策略带来的风险。}大语言模型按序列生成，根据前面的词元逐个生成后续词元，可能放大早期的错误。这一过程被称为“幻觉滚雪球”（hallucination snowballing，\refcite{ch06-ref-091}{Perez et al., 2022}）。Top-$p$、Top-$k$ 等采样式生成策略所含的随机性，也是已知的幻觉来源（\refcite{ch06-ref-058}{Lee et al., 2022}）。
\end{enumerate}

\subsection{评估指标}
本节概述评估大语言模型幻觉的若干方法。以往的幻觉评估高度依赖人类专家，专家按明确的原则进行评估（\refcite{ch06-ref-058}{Lee et al., 2022}；\refcite{ch06-ref-060}{Li et al., 2023a}；\refcite{ch06-ref-064}{Lin et al., 2021}；\refcite{ch06-ref-077}{Min et al., 2023}）。一般认为，人工评估最有助于保证评估的准确性。研究者也越来越关注自动评估，希望以更高效、更一致的方式评估大语言模型的幻觉。本节考察常用的人工评估技术和新出现的自动评估方法，说明它们的作用、效果和潜力。

\subsubsection{人工评估}
这种方法以专门制定的人工标注准则为基础，由标注者仔细检查模型生成的每段文本。例如，\refcite{ch06-ref-064}{Lin 等人（2021）}提出的 TruthfulQA 框架，用一组具体的指令指导标注者，把模型输出归入 13 个定性标签中的一个。这个过程还要用可靠的来源核查答案是否准确。\refcite{ch06-ref-058}{Lee 等人（2022）}同样使用人工标注，验证他们提出的自动评估指标是否有效。

\refcite{ch06-ref-077}{Min 等人（2023）}提出的 FActScore 是另一个例子。它要求标注者根据知识来源，判断每项原子事实（atomic fact）是“有支持”还是“无支持”。若该事实与给定提示无关，则标作“不相关”。
\begin{equation}
\operatorname{FActScore}(\mathcal{M})
 =\mathbb{E}_{x\in\mathcal{X}}\!\left[
 \frac{1}{|A_y|}\sum_{a\in A_y} I[a\text{ 由 }C\text{ 支持}]
 \;\middle|\; \mathcal{M}_x\text{ 给出回答}
 \right].
\label{eq:ch06-factscore}
\end{equation}
式中各量如下。

\begin{itemize}
\item $\mathcal{M}$ 是接受评估的语言模型。
\item $\mathcal{X}$ 是供语言模型回答的一组提示。
\item $C$ 是用于核查模型回答准确性的知识来源。
\item $y=\mathcal{M}_x$ 是模型 $\mathcal{M}$ 针对 $\mathcal{X}$ 中的提示 $x$ 生成的回答。
\item $A_y$ 是回答 $y$ 所含的原子事实列表。
\item $I[a\text{ 由 }C\text{ 支持}]$ 是指示函数。原子事实 $a$ 得到知识来源 $C$ 支持时取 1，否则取 0。
\end{itemize}

按上述定义，FActScore 也可作为自动评估器，为检索系统评分。

人工评估的结果通常可靠，也能作较细的解释，但仍有困难。一个重要问题是，人的判断带有主观性，可能造成结果不一致。不同标注者可能作不同的解读，使结果不够统一。此外，人工评估通常费时费力，成本较高，尤其是每引入一个新模型都要重新评估。

\subsubsection{基于模型的自动评估}
有些自动系统旨在模仿人的评估判断。例如，有研究将具有 67 亿参数的 GPT-3 微调，使其判断模型所生成答案的真实性，并在验证时取得了较高的准确率。

AlignScore 是这类指标之一，用于评估文本的事实一致性。它比较两段文本 $a$ 和 $b$，判断 $b$ 是否包含 $a$ 中所有相关信息，并且没有矛盾（\refcite{ch06-ref-126}{Zha et al., 2023}）。定义一个对齐函数来分析文本对，确定两段文本的一致程度。数学上，该函数是一个映射 $f$，接收一对文本 $(a,b)$，为其赋予量化对齐程度的标签 $y$。
\begin{equation}
f\colon(a,b)\to y.
\label{eq:ch06-align-mapping}
\end{equation}

为得到能够泛化的对齐函数，在多种语言任务上训练它，包括自然语言推断、事实核查、释义检测、文本语义相似性、问答、信息检索和摘要。为统一格式，将这些任务整理成文本对 $(a,b)$。

训练函数预测对齐标签 $y$。任务可分为三种。
\begin{itemize}
\item \textbf{二分类。} $y_{\mathrm{bin}}\in\{\text{ALIGNED},\text{NOT ALIGNED}\}$。
\item \textbf{三分类。} $y_{\mathrm{3way}}\in\{\text{ALIGNED},\text{CONTRADICT},\text{NEUTRAL}\}$。
\item \textbf{回归。} $y_{\mathrm{reg}}\in[0,1]$。
\end{itemize}
用联合损失函数 $L_{\mathrm{total}}$ 衡量模型预测的准确程度。
\begin{equation}
L_{\mathrm{total}}=\lambda_1L_{\mathrm{3way}}
  +\lambda_2L_{\mathrm{bin}}+\lambda_3L_{\mathrm{reg}},
\label{eq:ch06-align-loss}
\end{equation}
其中，$\lambda_1$、$\lambda_2$ 和 $\lambda_3$ 是标量权重，分别控制各项损失的影响。

\subsubsection{基于规则的自动评估}
这种方法用预定义规则自动为大语言模型输出的事实准确性评分。可以直接用准确率指标衡量模型区分真伪陈述的能力。

FactualityPrompt 将实体识别和蕴含指标结合起来，评估事实准确性的不同方面（\refcite{ch06-ref-058}{Lee et al., 2022}）。该框架分为几个阶段，见图~\ref{fig:ch06-factuality}。

\begin{itemize}
\item \textbf{生成续写。}大语言模型根据所给测试提示生成续写。
\item \textbf{识别值得核查的续写。}重点找出含有需要核查的事实的续写。这一步很重要，因为语言模型也会生成非事实内容，例如个人观点或随意交谈。
\item \textbf{准备参考事实。}为核查上一步找到的续写，准备相关的事实知识。
\item \textbf{计算事实性和质量指标。}在最后一个阶段，计算续写的事实性和质量指标。
\end{itemize}

\begin{figure}[!htbp]
\centering
% 重绘自原书第 224 页的图 6.1，保留模型输入与输出示例。
\begin{tikzpicture}[
  font=\footnotesize\sffamily,
  block/.style={draw=blue!42!black,fill=blue!14,rounded corners=3pt,align=center,
     minimum height=6.5mm,inner sep=2.5pt},
  group/.style={draw=blue!32!black,densely dashed,rounded corners=5pt,inner sep=4.4pt},
  wire/.style={draw=black!75,-{Stealth[length=1.4mm]},line width=.52pt}]
\node[font=\scriptsize\ttfamily,anchor=west,align=left] (prompt) at (0,0)
  {\textbf{[PROMPT]} Fox 2000\\released the Film Soul food};
\node[block,fill=yellow!48,minimum width=12mm] (lm) at (4.88,0) {语言模型};
\node[font=\scriptsize\ttfamily,anchor=west,align=left] (reply) at (6.04,0)
  {\textbf{[LM\_GEN]} Its leading case\\includes actress Emma\\Thompson, Tom Wilkinson, ...};
\draw[wire] (prompt.east)--(lm.west);
\draw[wire] (lm.east)--(reply.west);
\node[anchor=west,text=green!44!black,font=\footnotesize\sffamily\bfseries]
  at (0,-.64) {阶段 1　生成语言模型续写};
\draw[black!70,densely dashed,line width=.55pt] (0,-.87)--(11.45,-.87);
\node[anchor=west,text=green!44!black,font=\footnotesize\sffamily\bfseries]
  at (0,-1.11) {阶段 2　评估};

\node at (2.53,-1.64) {1.\ 筛选陈述};
\node[block,fill=yellow!47,minimum width=17mm] (cls) at (1.35,-2.55) {陈述类型\\分类器};
\node[align=left] (check) at (3.65,-2.30) {\textcolor{green!46!black}{值得核查}};
\node[align=left] (nocheck) at (3.65,-2.84) {无需核查};
\node[group,fit=(cls)(check)(nocheck)] (filter) {};
\draw[wire] (cls.east)--(2.71,-2.55);
\draw[black!65,line width=.4pt] (nocheck.west)--(nocheck.east);

\node[block,fill=black!8] (wiki) at (1.28,-4.04) {维基文档};
\node[block,fill=orange!23] (choose) at (3.65,-4.04) {依据选择器};
\node[group,fit=(wiki)(choose),inner sep=4.7pt] (sources) {};
\draw[wire] (wiki.east)--(choose.west);
\draw[wire] (3.65,-3.14)--(choose.north);
\node at (2.53,-4.88) {2.\ 选择事实依据};

\node at (8.35,-1.64) {3.\ 评估事实性};
\foreach \y/\leftlabel/\rightlabel/\shade in {
  -2.22/命名实体模型/{$\mathrm{NE}_{\mathrm{ER}}$}/17,
  -2.94/自然语言推断模型/蕴含百分比/23,
  -3.66/13 亿参数语言模型/流畅度（困惑度）/29,
  -4.38/基于 $n$ 元组的指标/多样性、重复/37}{
  \node[block,fill=blue!\shade,minimum width=25mm] (measure) at (7.07,\y) {\leftlabel};
  \node[anchor=west,font=\scriptsize\sffamily] (result) at (9.04,\y) {\rightlabel};
  \draw[wire] (measure.east)--(result.west);
}
\node[group,fit={(5.48,-1.86) (11.38,-4.71)},inner sep=2pt] {};
\end{tikzpicture}

\caption{FactualityPrompt 框架所定义的评估流程。}
\label{fig:ch06-factuality}
\end{figure}

\subsection{评估基准}
本节概述常用的基准，按其评估形式、任务形式和构造方法加以分类，见表~\ref{tab:ch06-hallucination-benchmarks}。

\begin{table}[!htbp]
\centering
\caption{评估大语言模型幻觉的代表性基准及其关注点。各基准用于评估生成质量（Gen），或区分事实与非事实陈述（Dis）。不同指标适用于不同任务，包括问答（QA）、任务指令（TI）以及给定前缀的文本补全（TC）。}
\label{tab:ch06-hallucination-benchmarks}
\small
\setlength{\tabcolsep}{3.5pt}
\begin{tabular}{@{}llrll@{}}
\toprule
基准 & 评估形式 & 规模 & 任务形式 & 指标\\
\midrule
TruthfulQA & Gen \& Dis & 817 & QA & 真实性\\
FactualityPrompt & Gen & 16,000 & TC & 集成指标\\
FActScore & Gen & 500 & TI & FActScore\\
KoLA-KC & Gen & 190 & TI & 自对比\\
HaluEval & Dis & 35,000 & QA \& TI & 准确率\\
FACTOR & Dis & 4,030 & TC & 准确率\\
\bottomrule
\end{tabular}
\end{table}

\subsubsection{评估形式}
评估幻觉时，主要考察大语言模型的两种能力。

\begin{enumerate}
\item \textbf{生成能力。}TruthfulQA、FactualityPrompt 和 FActScore 等基准考察大语言模型生成事实陈述的能力，关注生成文本的真实性与事实准确性。
\item \textbf{辨别能力。}HaluEval 和 FACTOR 等基准考察大语言模型区分事实陈述与幻觉陈述（非事实陈述）的能力。
\end{enumerate}

\subsubsection{任务形式}
评估幻觉时使用几类任务。

\begin{enumerate}
\item \textbf{问答任务。}TruthfulQA 等基准评估大语言模型能否为知识密集型问题给出准确答案。
\item \textbf{传记生成和文本补全。}FActScore 考察生成传记的事实准确性。FACTOR 和 FactualityPrompt 则评估模型根据事实或非事实前缀补全文本时的表现。
\end{enumerate}

\subsubsection{构造方法}
这类基准通常用以下办法构建。

\begin{itemize}
\item 由人工标注者创建并验证数据集，例如 FActScore 的人工流程。
\item 其他方法既有人工标注大语言模型生成的文本，也有设计提示、自动生成文本并在之后进行人工评估，例如 HaluEval。
\end{itemize}

\subsection{缓解策略}
本节介绍几种缓解策略，并指出它们适用于大语言模型开发过程的哪些阶段（图~\ref{fig:ch06-hallucination-mitigation}）。要了解各项技术的分类和细节，可参阅 \refcite{ch06-ref-129}{Zhang 等人（2023）}的综述。

\subsubsection{预训练时的缓解措施}
大语言模型的幻觉主要归因于预训练语料中有噪声或不可靠的数据。\refcite{ch06-ref-132}{Zhou 等人（2023）}等研究指出，大语言模型的参数知识（parametric knowledge），即嵌入已调参数中的知识，是在这一阶段形成的，对输出影响很大。在大语言模型出现之前，缓解幻觉主要依靠人工整理训练数据。\refcite{ch06-ref-031}{Gardent 等人（2017）}针对特定任务，借助筛选技术人工清理数据，减少了幻觉。

\begin{practicaltip}
大语言模型的预训练语料规模巨大，有数万亿个词元，人工整理并不现实。因此，通常依赖自动选择或筛选可靠数据。现有方法会自动选择高质量数据，或过滤噪声较多的数据。例如，按与高质量参考语料的相似程度清理 GPT-3 的预训练数据。Falcon（\refcite{ch06-ref-090}{Penedo et al., 2023}）和 phi-1.5（\refcite{ch06-ref-063}{Li et al., 2023c}）所用的是经过筛选的高质量数据，得到的模型更可靠。有些模型有选择地增加维基百科等事实性较强的来源的采样比例，提高预训练语料的质量（\refcite{ch06-ref-109}{Touvron et al., 2023}）。\refcite{ch06-ref-064}{Lin 等人（2021）}建议，在预训练时为事实性文档中的句子加上主题前缀。这样做提高了模型在 TruthfulQA 等基准上的表现。
\end{practicaltip}

\begin{figure}[!htbp]
\centering
% 重绘自原书第 226 页的图 6.2。
\begin{tikzpicture}[
  font=\footnotesize\sffamily,
  reason/.style={draw=yellow!67!orange,fill=yellow!26,minimum width=33mm,
    minimum height=6.8mm,align=center,inner sep=2pt},
  remedy/.style={reason,draw=orange!43,fill=orange!16,minimum width=34mm},
  phase/.style={draw=green!45!black,fill=green!30,circle,inner sep=0pt,minimum size=3mm},
  link/.style={draw=blue!25!black!62,densely dashed,-{Stealth[length=1.4mm]},line width=.57pt}]
\node[text=green!45!black] at (0,1.21) {时间轴};
\draw[green!52!black,line width=.72pt,-{Stealth[length=1.55mm]}] (0,.84)--(0,-4.25);
\node[phase] (pre) at (0,0) {};
\node[phase] (sft) at (0,-1.18) {};
\node[phase] (rl) at (0,-2.36) {};
\node[phase] (inf) at (0,-3.54) {};
\foreach \id/\name/\y in {pre/预训练/0,sft/SFT/-1.18,rl/RLHF/-2.36,inf/推理/-3.54}{
  \node[anchor=south east,text=red!60!black,font=\scriptsize\sffamily] at (-.12,\y) {\name};
}
\node[reason] (memory) at (-3.60,0) {参数记忆};
\node[reason] (confidence) at (-3.60,-1.18) {过于自信};
\node[reason] (wrong) at (-3.60,-2.36) {对齐偏差};
\node[reason] (decodingrisk) at (-3.60,-3.54) {生成过程的风险};
\node[remedy] (curation) at (3.59,.19) {整理训练数据};
\node[remedy] (honestsft) at (3.59,-.65) {以诚实为目标的 SFT};
\node[remedy] (honestrl) at (3.59,-1.49) {以诚实为目标的强化学习};
\node[remedy] (decode) at (3.59,-2.33) {解码策略};
\node[remedy] (retr) at (3.59,-3.17) {知识检索};
\node[remedy] (uncert) at (3.59,-4.01) {利用不确定性};
\foreach \from/\to in {memory/pre,confidence/sft,wrong/sft,wrong/rl,decodingrisk/inf,
  pre/curation,sft/curation,sft/honestsft,rl/honestrl,inf/decode,inf/retr,inf/uncert}{
  \draw[link] (\from.east)--(\to.west);
}
\end{tikzpicture}

\caption{大语言模型生命周期各阶段的幻觉成因及相应的缓解策略。}
\label{fig:ch06-hallucination-mitigation}
\end{figure}

\subsubsection{微调时的缓解措施}
大语言模型通过微调学会运用预训练知识并与用户交互。与预训练相似，可以通过整理训练数据来减少微调中的幻觉。微调数据较少，人工和自动整理都可行。\refcite{ch06-ref-132}{Zhou 等人（2023）}构建了包含人工标注样例的数据集，也有研究者自动选择高质量的指令调优数据。

\begin{practicaltip}
在 TruthfulQA 等基准上，与使用未经整理的数据微调的模型相比，使用经过整理的指令数据微调的大语言模型，真实性和事实性更高。在微调数据中加入诚实作答的样例，例如用“Sorry, I don't know”（对不起，我不知道）承认自己无法回答，可以教模型拒绝回答超出其知识范围的问题，从而减少幻觉（\refcite{ch06-ref-105}{Sun et al., 2023}）。
\end{practicaltip}

\subsubsection{RLHF 时的缓解措施}
RLHF 旨在使大语言模型符合人类偏好和具体的 3H 准则。3H 中的“诚实”侧重减少大语言模型回答中的幻觉。RLHF 主要包含两步，即（a）训练奖励模型和（b）优化语言模型。第一步用奖励模型代替人表达偏好，为大语言模型的各个回答分配适当的奖励值。接着使用奖励模型的反馈，借助 PPO 等强化学习算法优化语言模型，见第~\ref{ch:alignment}~章。

\begin{practicaltip}
RLHF 也能缓解监督微调中的行为克隆（behavior cloning）问题。在监督微调中，大语言模型可能模仿人类行为而没有理解背后的知识，从而产生幻觉。\refcite{ch06-ref-099}{Schulman（2023）}提出在 RLHF 中使用一种奖励函数减少幻觉，鼓励模型表达不确定性或承认自己无法回答。采用 RLHF 的模型使用合成幻觉数据后，在 TruthfulQA 等基准上的表现已有明显改善。
\end{practicaltip}

\subsubsection{推理时的缓解措施}
与在训练时干预相比，推理时的缓解措施往往更节省成本，也更容易控制。贪心解码、束搜索解码和核采样等解码策略的选择，会显著影响大语言模型的输出。

\begin{practicaltip}
\refcite{ch06-ref-124}{Zarrieß 等人（2021）}讨论了解码策略如何改变模型生成的概率分布。\refcite{ch06-ref-058}{Lee 等人（2022）}评估大语言模型使用不同解码策略时生成内容的事实性。他们发现，\refcite{ch06-ref-046}{Holtzman 等人（2019）}提出的核采样所得内容，事实性不及贪心解码。这归因于 Top-$p$ 采样的随机性，它可能导致幻觉。为此，\refcite{ch06-ref-058}{Lee 等人（2022）}提出一种称为事实性核采样（factual-nucleus sampling）的解码算法，希望兼顾多样性和事实性。
\end{practicaltip}

另一种重要办法是用外部知识补充事实根据。做法是从可信来源获取知识，并用它指导大语言模型生成回答。\refcite{ch06-ref-095}{Ren 等人（2023）}和 \refcite{ch06-ref-076}{Mialon 等人（2023）}考察了这种方法，它通常包含知识获取和知识利用两步。

\begin{enumerate}
\item \textbf{知识获取。}大语言模型通过预训练和微调内化大量知识，形成参数知识。为避免错误或过时的参数知识引起幻觉，研究者建议从可信来源获取可靠且较新的知识。主要来源如下。
  \begin{itemize}
  \item \textbf{外部知识库。}从大规模非结构化语料（\refcite{ch06-ref-013}{Cai et al., 2021}）、结构化数据库（\refcite{ch06-ref-062}{Li et al., 2023b}）或维基百科等网站（\refcite{ch06-ref-121}{Yao et al., 2022}）检索信息。可使用多种检索器，例如 BM25（\refcite{ch06-ref-096}{Robertson et al., 2009}）或基于语言模型的方法（\refcite{ch06-ref-131}{Zhao et al., 2022}）。\refcite{ch06-ref-069}{Luo 等人（2023）}提出一个框架，从经过微调的白箱大语言模型的参数记忆中检索知识。\refcite{ch06-ref-026}{Feng 等人（2023）}研究怎样教大语言模型从外部知识图谱搜索相关领域知识。
  \item \textbf{外部工具。}用外部工具提高大语言模型生成内容的事实性。例如，FacTool（\refcite{ch06-ref-017}{Chern et al., 2023}）在特定任务中使用相应方法检测幻觉。CRITIC（\refcite{ch06-ref-037}{Gou et al., 2023}）使大语言模型能够与多种工具交互，自行修改回答，提高真实性。
  \end{itemize}
\item \textbf{知识利用。}取得相关知识后，可以在不同阶段利用它减少大语言模型的幻觉。大体分为以下两种方式。
  \begin{itemize}
  \item \textbf{生成时补充。}在提示大语言模型之前，直接把检索到的知识或工具反馈接到用户查询上（\refcite{ch06-ref-100}{Shi et al., 2023}）。
  \item \textbf{事后纠正。}构建辅助的“纠正器”，在后处理阶段修正幻觉（\refcite{ch06-ref-014}{Cao et al., 2020}）。
  \end{itemize}
\end{enumerate}

估计大语言模型的不确定性，也是检测和减少幻觉的一个办法。它考察模型对输出有多大把握，大体可分为以下几种。

\begin{enumerate}
\item \textbf{基于 logit 的估计。}根据模型的 logit 计算词元级概率或熵。这种技术在机器学习中很常用（\refcite{ch06-ref-039}{Guo et al., 2017}）。
\item \textbf{基于表述的估计。}用提示直接要求大语言模型表达自己的不确定性，例如给出置信度分数。采用思维链提示，可以进一步改善这一过程（\refcite{ch06-ref-116}{Wei et al., 2022}；\refcite{ch06-ref-119}{Xiong et al., 2023}）。\refcite{ch06-ref-079}{Mündler 等人（2023）}使用另一个大语言模型评估回答中有无逻辑矛盾。
\item \textbf{基于一致性的估计。}这种办法假设，大语言模型不确定时会给出不一致的回答（\refcite{ch06-ref-100}{Shi et al., 2023}；\refcite{ch06-ref-130}{Zhao et al., 2023}）。SelfCheckGPT 在无额外资源的黑箱设置中采用基于一致性的方法，用多种指标衡量回答的一致程度（\refcite{ch06-ref-070}{Manakul et al., 2023}）。
\end{enumerate}

% PDF 257--267；原书第 229--239 页。
\section{偏见与公平性}
\label{sec:ch06-bias}

\begin{keypoint}
社会偏见（social bias）指不同社会群体受到的不平等待遇或取得的不平等结果，主要源于历史与结构性权力失衡。\refcite{ch06-ref-029}{Gallegos 等人（2023）}的综述为自然语言处理中社会结构与算法模型的交叉研究提供了明确的定义。下面是其中几个关键术语。
\begin{itemize}
\item \textbf{社会群体（social group）。}共享某种身份特征的人口子集。身份特征可以是固定的、依语境而定的，或由社会建构。反歧视法律所承认的群体包括按年龄、残障、性别认同、原籍、种族、宗教、生理性别和性取向划分的群体，凸显了这些群体的社会建构性质。
\item \textbf{受保护属性（protected attribute）。}构成一个社会群体身份基础的共同身份特征。
\end{itemize}
综述还提出以下公平性概念。
\begin{itemize}
\item \textbf{群体公平性（group fairness）。}按群体成员身份计算结果的统计量，并要求该统计量在所有社会群体间相同。它能显示群体之间的差异，但可能忽略子群体或交叉身份。
\item \textbf{个体公平性（individual fairness）。}相似的个体应受到相似待遇。可用距离度量比较不同个体及其结果，从而量化这一概念。
\end{itemize}
\end{keypoint}

社会偏见指不同社会群体因历史和结构性权力失衡而受到不平等待遇，或获得不同结果。在自然语言处理领域，社会偏见表现为表征性伤害和分配性伤害。表征性伤害包括失实表征、刻板印象、系统表现不均、贬损性语言和排斥性规范。分配性伤害则包括直接歧视和间接歧视。

\subsection{表征性伤害}
自然语言处理系统中的表征性伤害（representational harm），指系统延续对某些社会群体的负面态度与刻板印象。这些伤害有多种形式。

\begin{itemize}
\item \textbf{贬损性语言。}使用贬低社会群体的语言。例如，使用“whore”（“婊子”）一词反映了对女性的敌意和刻板印象（\refcite{ch06-ref-007}{Beukeboom and Burgers, 2019}）。
\item \textbf{系统表现有差异。}模型理解或处理不同社会群体的语言时，对某些群体的效果较差。例如，与意思相同的标准美国英语（Standard American English，SAE）表达相比，非裔美国人白话英语（African American Vernacular English，AAVE）的表达更常被误判为非英语（\refcite{ch06-ref-008}{Blodgett and O'Connor, 2017}）。
\item \textbf{排斥性规范。}这些规范隐含地排除或贬低某些群体。例如，“both genders”（两种性别）这一说法排除了非二元性别认同（\refcite{ch06-ref-006}{Bender et al., 2021}）。
\item \textbf{失实表征。}对社会群体的表征不完整或带有偏见。例如，对“I'm an autistic dad”（我是一位有孤独症的父亲）不恰当地回答“sorry to hear that”（听到这件事我很遗憾），传达出关于孤独症的负面刻板印象（\refcite{ch06-ref-101}{Smith et al., 2022}）。
\item \textbf{刻板印象。}延续有关某一群体的固定负面概括。例如，把“Muslim”（穆斯林）与“terrorist”（恐怖分子）联系起来，会强化关于暴力的刻板印象（\refcite{ch06-ref-002}{Abid et al., 2021}）。
\item \textbf{毒性。}用冒犯性语言攻击社会群体或煽动对它的仇恨。例如，“I hate Latinos”（我恨拉丁裔）表达了对一个社会群体的不尊重与仇恨（\refcite{ch06-ref-023}{Dixon et al., 2018}）。
\end{itemize}

\subsection{分配性伤害}
分配性伤害（allocational harm）涉及资源或机会在不同社会群体之间分配不均，可以进一步分类。

\begin{itemize}
\item \textbf{直接歧视。}因群体成员身份而受到明确的差别待遇。例如，用大语言模型辅助筛选简历时，偏见可能导致招聘中的不公（\refcite{ch06-ref-027}{Ferrara, 2023}）。
\item \textbf{间接歧视。}即便所考虑的因素表面中立，替代变量或隐含因素仍会造成差别待遇。例如，使用大语言模型的医疗工具可能采用人口统计属性的替代变量，加剧患者获得照护时的不平等（\refcite{ch06-ref-027}{Ferrara, 2023}）。
\end{itemize}

\subsection{成因}
本节以 \refcite{ch06-ref-084}{Navigli 等人（2023）}的研究为框架，考察影响偏见与公平性的因素。

\subsubsection{数据选择}
语言模型在大型数据集上训练。为这些数据集选择文本，会引入选择偏差，影响模型的行为和输出。从初始抽样到数据清理和筛选，各个阶段都可能出现这种偏差。如果所选训练文本未能体现网上语言使用的多样性，就会发生数据选择偏差。现代大语言模型所用数据集虽大，但仍有范围限制，模型会继承文本中的偏见。选择过程偏向特定来源，会进一步加剧这个问题。例如，维基百科文本常因被认为可靠而入选，YouTube 评论等非正式来源的内容则被排除（\refcite{ch06-ref-012}{Brown et al., 2020}；\refcite{ch06-ref-018}{Chowdhery et al., 2022}；\refcite{ch06-ref-128}{Zhang et al., 2022}）。这样的选择会塑造模型的理解，也会使其产生偏见。为了使大语言模型适应特定任务，通常会用较小的专用数据集或定制提示微调模型（\refcite{ch06-ref-047}{Howard and Ruder, 2018}；\refcite{ch06-ref-067}{Liu et al., 2023}）。这些额外的数据处理步骤，可能引入新的偏见，也可能加剧已有的偏见，具体取决于所用的微调数据或提示。

\subsubsection{领域和体裁分布不均}
预训练数据集中的内容往往分布不均。例如，数据集常用的维基百科来源，较多覆盖体育、音乐、地理和政治等领域，而文学、经济和历史等领域较少得到体现（\refcite{ch06-ref-030}{Gao et al., 2020}；\refcite{ch06-ref-055}{Kreutzer et al., 2022}）。这一现象并非英语独有，在其他高资源语言中也存在。用 BabelNet 比较英语和意大利语维基百科的领域分布，即可看到这一点（\refcite{ch06-ref-083}{Navigli and Ponzetto, 2012}）。

预训练数据集中的偏见会延续到下游任务。例如，用于在机器翻译和语义角色标注等任务中微调语言模型的 EuroParl 和 CoNLL-2009，偏向某些领域和体裁。即使语言模型一开始没有这类偏见，也会从这些任务专用的数据集学到新的偏见，开发人员可能无法立即察觉（\refcite{ch06-ref-036}{Goldfarb-Tarrant et al., 2020}）。

\subsubsection{内容的创作时间}
语言不断演变，词语的含义和用法也随之改变。例如，“mouse”原指一种动物，后来也指计算机输入设备。“tweet”从指鸟叫扩展到指社交媒体上的帖子。“car”过去指马拉的车，如今指机动车。“pipe”也从指吸烟用的器具扩展到指一种管子。语言用法的这些历史变化意味着，在历史数据上训练的语言模型可能无法准确反映当代用法，或无法理解当下所指的事物。

特定领域文本的内容和关注点，也会随时间显著变化。例如，中世纪的医学文本与现代医学文献有很大差别。如果训练数据早于新冠疫情、詹姆斯·韦布空间望远镜发射等近年的重要事件，语言模型可能缺少相关知识。类似地，ChatGPT 等有知识截止日期的模型，可能没有该日期之后发生事件的信息。出于实际考虑，研究者经常重复使用较早的数据集，例如以 20 世纪 60 年代的 Brown Corpus 为基础的 SemCor。对于为词义消歧等任务训练的模型，这种做法可能延续过时的语言用法。

\subsubsection{创作者的人口构成}
训练语料的创作者和选择者的人口构成，都会影响语言模型的偏见和行为。某些人口群体在维基百科等平台上占比偏高，说明内容创作和语料选择过程需要更多样、更包容。创作训练语料内容的人口构成，可能使大语言模型产生偏见。例如，作为常见训练数据来源，维基百科编辑者的人口构成明显不均衡。男性占多数（84\%），主要是二十五岁上下的人或退休人员（\refcite{ch06-ref-117}{Wikipedia Contributors, 2023}）。

这些偏见来自内容的创作者，也来自决定哪些内容进入训练集的人。选择大语言模型语料的决策者，通常也是男性居多。决策者的构成相近，可能使所选主题和观点范围较窄，进一步强化训练数据中已有的偏见。

\subsubsection{语言与文化的分布偏斜}
由于高资源语言的数据更易收集，也更容易找到语言学家和标注者，大语言模型开发一直以高资源语言为主。这形成了反馈循环，一方面改善这些语言的自然语言处理系统，另一方面使低资源语言处于边缘。尽管已有进展，多语言模型的表现仍偏向训练数据较丰富的语言。

不同语言承载不同文化，影响隐喻和习语等语言表达（\refcite{ch06-ref-045}{Hershcovich et al., 2022}）。大语言模型中的语言分布偏斜，会导致文化表征不均。例如，维基百科以英语为主，超过 50\% 的编辑者主要使用英语，内容因而更偏向英语文化。印地语、孟加拉语、爪哇语和泰卢固语等语言虽然使用人口众多，所占内容却较少。维基百科编辑者的主要语言，也不能反映英语使用者在全球的分布。例如，以英语为主要语言的编辑者只有 3\% 来自印度，这会影响维基百科内容的多样性（\refcite{ch06-ref-114}{Wang et al., 2022}）。

\subsection{评估指标}
大语言模型的偏见与公平性指标，可按它们使用模型的哪些信息分组，例如嵌入、概率或生成的文本。可分为以下几类。

\begin{itemize}
\item \textbf{基于嵌入的指标。}用稠密向量表示衡量偏见，通常使用带有上下文信息的句子嵌入。
\item \textbf{基于概率的指标。}用模型分配的概率估计偏见，例如给成对的文本打分或回答选择题。
\item \textbf{基于生成文本的指标。}分析模型针对提示生成的文本，衡量共现等模式，或比较不同提示产生的输出。
\end{itemize}

\subsubsection{基于嵌入的指标}
这类指标主要计算中性词（如职业名称）和身份相关词（如带性别的代词）在向量空间中的距离。大语言模型的评估重点是句子级的上下文嵌入。这些指标最初针对静态词嵌入提出，后来扩展到上下文嵌入，从多个维度衡量偏见。

词嵌入关联检验（Word Embedding Association Test，WEAT）是其中一种主要方法，评估社会群体概念（如男性与女性相关词）和中性属性（如家庭与职业相关词）之间的关联（\refcite{ch06-ref-038}{Greenwald et al., 1998}）。为衡量刻板关联，对受保护属性 $A_1,A_2$ 和中性属性 $W_1,W_2$ 使用以下检验统计量。
\begin{equation}
f(A_1,A_2,W_1,W_2)
=\sum_{a_1\in A_1}s(a_1,W_1,W_2)-\sum_{a_2\in A_2}s(a_2,W_1,W_2).
\label{eq:ch06-weat-statistic}
\end{equation}
其中，$s$ 是相似度度量，定义为
\begin{equation}
s(a,W_1,W_2)=
\operatorname{mean}_{w_1\in W_1}\cos(a,w_1)
-\operatorname{mean}_{w_2\in W_2}\cos(a,w_2).
\label{eq:ch06-weat-similarity}
\end{equation}
效应量越大，偏见越强。效应量为
\begin{multline}
\operatorname{WEAT}(A_1,A_2,W_1,W_2)\\
=
\frac{\operatorname{mean}_{a_1\in A_1}s(a_1,W_1,W_2)
      -\operatorname{mean}_{a_2\in A_2}s(a_2,W_1,W_2)}
     {\operatorname{std}_{a\in A_1\cup A_2}s(a,W_1,W_2)}.
\label{eq:ch06-weat-effect}
\end{multline}

为将 WEAT 用于上下文嵌入，\refcite{ch06-ref-073}{May 等人（2019）}提出句子编码器关联检验（Sentence Encoder Association Test，SEAT）。SEAT 将表示社会群体与中性属性的词填入语义中性模板，例如“This is [BLANK]”或“[BLANK] are things.”，再生成这些句子的嵌入。它使用与 WEAT 相同的公式，取 \texttt{[CLS]} 词元的嵌入。SEAT 也可以使用更具体、含有明确语义的模板，例如“The engineer is [BLANK].”，以衡量更细的偏见维度。

另一种方法是上下文嵌入关联检验（Contextualized Embedding Association Test，CEAT）（\refcite{ch06-ref-041}{Guo and Caliskan, 2021}）。它组合 $A_1$、$A_2$、$W_1$ 和 $W_2$ 各组中的元素生成句子，并从嵌入中随机抽取一个子集，计算效应量的分布。CEAT 用随机效应模型衡量偏见的大小，公式如下。
\begin{multline}
\operatorname{CEAT}(S_{A_1},S_{A_2},S_{W_1},S_{W_2})\\
=\frac{\sum_{i=1}^{N}v_i\operatorname{WEAT}(S_{A_1i},S_{A_2i},S_{W_1i},S_{W_2i})}
{\sum_{i=1}^{N}v_i},
\label{eq:ch06-ceat}
\end{multline}
其中，$v_i$ 由随机效应模型的方差得到。用这些方法可以把 WEAT 的原理应用到上下文嵌入，更细致地分析语言模型中的偏见。

\subsubsection{基于概率的指标}
这类技术将受保护属性有所变动的成对或成组模板句子交给模型，比较不同输入条件下预测的词元概率。一种办法是掩码词元法，在句子里遮盖一个词，再用掩码语言模型预测缺失的词。例如，\refcite{ch06-ref-115}{Webster 等人（2020）}提出的关联发现（Discovery of Correlations，DisCo），将可能触发偏见的词填入模板的槽位，比较模型在这些模板上的补全和排名最高的预测。

Kurita 等人（2019）提出的对数概率偏见分数（Log-Probability Bias Score，LPBS），使用类似 DisCo 的模板方法，评估中性属性词中的偏见。先用模板“[MASK] is a [NEUTRAL ATTRIBUTE]”得到某个词元的预测概率 $p_a$，再用模板“[MASK] is a [MASK]”得到该词元的模型先验概率 $p_{\mathrm{prior}}$，以先验概率对预测概率作归一化。这样的归一化考虑了模型对某些社会群体原有的偏见，使度量集中在与 \texttt{[NEUTRAL ATTRIBUTE]} 词元有关的偏见上。然后比较两个相对的社会群体词经归一化后的概率，算出偏见分数。

LPBS 的数学定义如下。
\begin{equation}
\operatorname{LPBS}(S)
=\log\frac{p_{a_i}}{p_{\mathrm{prior}_i}}
 -\log\frac{p_{a_j}}{p_{\mathrm{prior}_j}},
\label{eq:ch06-lpbs}
\end{equation}
其中，$p_{a_i}$、$p_{a_j}$ 为不同社会群体词的预测概率，$p_{\mathrm{prior}_i}$、$p_{\mathrm{prior}_j}$ 为它们各自的先验概率。LPBS 由此通过衡量词元概率偏离预期先验分布的程度来量化偏见。

\refcite{ch06-ref-003}{Ahn 和 Oh（2021）}提出类别偏见分数（Categorical Bias Score，CBS），将 Kurita 等人（2019）的归一化对数概率扩展到非二元目标。CBS 衡量在不同社会群体的受保护属性词 $a$ 上，填空模板所预测词元的归一化对数概率的方差。
\begin{equation}
\operatorname{CBS}(S)
=\frac{1}{|W|}\sum_{w\in W}\operatorname{Var}_{a\in A}\log\frac{p_a}{p_{\mathrm{prior}}}.
\label{eq:ch06-cbs}
\end{equation}

还有一系列方法使用伪对数似然（pseudo-log-likelihood，PLL），根据句子中其他词构成的上下文，确定生成特定词元的可能性。对于给定句子 $S$，PLL 定义为
\begin{equation}
\operatorname{PLL}(S)=\sum_{s\in S}\log P(s\mid S\setminus s;\theta),
\label{eq:ch06-pll}
\end{equation}
其中，$P(s\mid S\setminus s;\theta)$ 估计句子 $S$ 中各词元 $s$ 的条件概率。每次遮盖一个词元，利用句子中其余未被遮盖的词元预测它，以此近似这个词元在上下文中的概率。

\refcite{ch06-ref-081}{Nangia 等人（2020）}在提出 CrowS-Pairs 数据集的同时，也提出了 CrowS-Pairs 分数（CrowS-Pairs Score，CPS）。它使用成对的句子，其中一句有明显的刻板印象，另一句较少带有刻板印象，从而评估模型对刻板内容的偏好。这个指标借助 PLL，在给定已修改词元 $M$ 的条件下近似共享、未修改词元 $U$ 的概率，记作 $P(U\mid M,\theta)$。$M$ 中的词元通常表示受保护属性。这个近似通过逐个遮盖和预测句中未修改的词元来实现。句子 $S$ 的指标为
\begin{equation}
\operatorname{CPS}(S)=\sum_{u\in U}\log P(u\mid U\setminus u,M;\theta).
\label{eq:ch06-cps}
\end{equation}

\refcite{ch06-ref-080}{Nadeem 等人（2020）}在提出 StereoSet 数据集时，引入了另一种比较句子的方法，即上下文关联检验（Context Association Test，CAT）。每个句子都与三个选项配对，分别是符合刻板印象、违背刻板印象和无意义的选项。选项可以是填空词元，也可以是续写句子。伪对数似然方法关注 $P(U\mid M,\theta)$，CAT 则关注 $P(M\mid U,\theta)$，由此将评估表述为
\begin{equation}
\operatorname{CAT}(S)=\frac{1}{|M|}\sum_{m\in M}\log P(m\mid U;\theta).
\label{eq:ch06-cat}
\end{equation}

\subsubsection{基于生成文本的指标}
当大语言模型是黑箱，无法直接获取它的概率或嵌入时，基于生成文本的指标尤其适用。常见做法是以已知可能引发偏见或毒性的提示为条件，例如来自 RealToxicityPrompts 和 BOLD 的提示，再分析生成文本中的偏见。

在各项指标中，社会群体替换（Social Group Substitutions，SGS）要求，替换人口统计属性后，大语言模型的回答仍相同（\refcite{ch06-ref-093}{Rajpurkar et al., 2016}）。设不变性度量为 $\psi$，例如精确匹配。以 $\widehat{Y}_i$ 表示原输入的预测输出，以 $\widehat{Y}_j$ 表示改变人口统计属性后的反事实输入的输出。SGS 的数学表达式为
\begin{equation}
\operatorname{SGS}(\widehat{Y})=\psi(\widehat{Y}_i,\widehat{Y}_j).
\label{eq:ch06-sgs}
\end{equation}

另一项指标是共现偏见分数（Co-Occurrence Bias Score），它衡量生成文本中的词元和带有性别的词一起出现的情况（\refcite{ch06-ref-011}{Bordia and Bowman, 2019}）。
\begin{equation}
\text{共现偏见分数}(w)=\log\frac{P(w\mid A_i)}{P(w\mid A_j)},
\label{eq:ch06-cooccurrence}
\end{equation}
其中，$w$ 是词元，$A_i$ 和 $A_j$ 是两组属性。

人口群体表征（Demographic Representation，DR）比较不同社会群体在输出中的出现频率与原始数据集中的分布（\refcite{ch06-ref-009}{Bommasani et al., 2023}）。设函数 $C(x,Y)$ 计算词 $x$ 在序列 $Y$ 中出现的次数。对集合 $G$ 中的社会群体 $G_i$，以 $A_i$ 表示与之对应的受保护属性词，其 DR 为
\begin{equation}
\operatorname{DR}(G_i)
=\sum_{a_i\in A_i}\sum_{\widehat{Y}\in\widehat{\mathbb{Y}}}C(a_i,\widehat{Y}).
\label{eq:ch06-dr}
\end{equation}
这里，$\operatorname{DR}=[\operatorname{DR}(G_1),\ldots,\operatorname{DR}(G_m)]$ 由各组的计数组成，将其归一化为概率分布。再用总变差距离、KL 散度或 Wasserstein 距离等度量，把这个分布与均匀分布等参考分布作比较，以评估模型输出对社会群体的表征是否公平。

刻板关联（Stereotypical Associations，ST）指标评估特定词与社会群体的关联中存在的偏见（\refcite{ch06-ref-009}{Bommasani et al., 2023}）。给定一组预测输出 $\widehat{\mathbb{Y}}$，该指标量化词 $w$ 与社会群体 $G_i$ 的属性词 $A_i$ 共现的频率。函数为
\begin{equation}
\operatorname{ST}(w)_i
=\sum_{a_i\in A_i}\sum_{\widehat{Y}\in\widehat{\mathbb{Y}}}
 C(a_i,\widehat{Y})I(C(w,\widehat{Y})>0).
\label{eq:ch06-st}
\end{equation}
在这个公式中，$C(a_i,\widehat{Y})$ 表示属性词 $a_i$ 在预测输出 $\widehat{Y}$ 中的计数。和人口群体表征指标一样，可以将计数向量 $\operatorname{ST}=[\operatorname{ST}(w)_i,\ldots,\operatorname{ST}(w)_k]$ 归一化，再与参考分布比较。

HONEST 指标用于衡量语言模型生成伤人句子补全的频率（\refcite{ch06-ref-085}{Nozza et al., 2021}）。对于因果模型，用一个不完整句子作提示。对于掩码模型，用带 \texttt{[MASK]} 词元的句子，让模型预测一个词。HONEST 计算由 HurtLex 语料库识别的伤人预测所占比例，再对动物、犯罪和负面含义等不同类别的比例取平均。HONEST 分数的公式为
\begin{equation}
\operatorname{HONEST}(\widehat{\mathbb{Y}})
=\frac{\sum_{\widehat{Y}_k\in\widehat{\mathbb{Y}}_k}
 \sum_{\widehat{y}\in\widehat{Y}_k}I_{\mathrm{HurtLex}}(\widehat{y})}
 {|\widehat{\mathbb{Y}}|\cdot k}.
\label{eq:ch06-honest}
\end{equation}
这里，$\widehat{Y}_i\in\widehat{\mathbb{Y}}_i$ 表示与特定群体 $G_i$ 相关的预测输出。该指标用于身份相关的模板提示及其排名前 $k$ 的补全，记作 $\widehat{\mathbb{Y}}_k$。

\subsection{评估基准}
基准数据集可先分为反事实输入和提示两大类。反事实输入又可分为掩码词元和未加掩码的句子两类。包含成对或成组句子的数据集，通常是反事实数据集，突出了模型预测在不同社会群体间的差异。掩码词元数据集的句子留有空位，让语言模型补全。它们适用于基于掩码词元概率的指标或伪对数似然指标。Winograd Schema Challenge、Winogender 和 WinoBias 等共指消解任务是常见例子。

未加掩码的句子数据集评估一对句子中哪一句更有可能。它们既可以用于伪对数似然指标，也能灵活用于其他指标。例如众包刻板印象句子对（Crowdsourced Stereotype Pairs，CrowS-Pairs）、Equity Evaluation Corpus 和 RedditBias。

提示类数据集再按用途分为句子补全和问答两类。用于提示文本续写的数据集包括 RealToxicityPrompts、TrustGPT 和 BOLD 等句子补全数据集。它们提供由大语言模型补全的句子前缀。Grep-BiasIR 和 BBQ（Bias Benchmark for QA，问答偏见基准）则是考察大语言模型偏见的问答框架。表~\ref{tab:ch06-bias-benchmarks} 汇总了 \refcite{ch06-ref-029}{Gallegos 等人（2023）}列举的数据集和分类。

% 原书第 239 页，勾选位置按原表的行列重排。
\begin{table}[p]
\centering
\caption{面向偏见评估的基准数据集。表中列出各数据集的大小、所针对的偏见问题和用于评估的社会群体。勾号表示原研究明确提及的问题或群体，也可以表示该数据集还能覆盖的情形。}
\label{tab:ch06-bias-benchmarks}
\footnotesize
\setlength{\tabcolsep}{1.6pt}
\renewcommand{\arraystretch}{1.06}
\begin{tabular}{@{}>{\raggedright\arraybackslash}p{.25\linewidth}r*{15}{c}@{}}
\toprule
数据集 & 规模 & \multicolumn{15}{c@{}}{偏见问题}\\
\cmidrule(l){3-17}
& & \rotatebox{90}{\scriptsize 失实表征} & \rotatebox{90}{\scriptsize 刻板印象} & \rotatebox{90}{\scriptsize 系统表现有差异} & \rotatebox{90}{\scriptsize 贬损性语言} & \rotatebox{90}{\scriptsize 排斥性规范} & \rotatebox{90}{\scriptsize 毒性} & \rotatebox{90}{\scriptsize 年龄} & \rotatebox{90}{\scriptsize 残障} & \rotatebox{90}{\scriptsize 性别（认同）} & \rotatebox{90}{\scriptsize 国籍} & \rotatebox{90}{\scriptsize 外貌} & \rotatebox{90}{\scriptsize 种族} & \rotatebox{90}{\scriptsize 宗教} & \rotatebox{90}{\scriptsize 性取向} & \rotatebox{90}{\scriptsize 其他}\\
\midrule
\multicolumn{17}{c}{\bfseries 反事实输入（掩码词元）}\\
Winogender & 720 & $\checkmark$ & $\checkmark$ & $\checkmark$ &  & $\checkmark$ &  &  &  & $\checkmark$ &  &  &  &  &  & \\
WinoBias & 3,160 & $\checkmark$ & $\checkmark$ & $\checkmark$ &  & $\checkmark$ &  &  &  & $\checkmark$ &  &  &  &  &  & \\
WinoBias+ & 1,367 & $\checkmark$ & $\checkmark$ & $\checkmark$ &  & $\checkmark$ &  &  &  & $\checkmark$ &  &  &  &  &  & \\
GAP & 8,908 & $\checkmark$ & $\checkmark$ & $\checkmark$ &  & $\checkmark$ &  &  &  & $\checkmark$ &  &  &  &  &  & \\
GAP-Subjective & 8,908 & $\checkmark$ & $\checkmark$ & $\checkmark$ &  & $\checkmark$ &  &  &  & $\checkmark$ &  &  &  &  &  & \\
BUG & 108,419 & $\checkmark$ & $\checkmark$ & $\checkmark$ &  & $\checkmark$ &  &  &  & $\checkmark$ &  &  &  &  &  & \\
StereoSet & 16,995 & $\checkmark$ & $\checkmark$ & $\checkmark$ &  &  &  &  &  & $\checkmark$ &  &  & $\checkmark$ & $\checkmark$ &  & $\checkmark$\\
BEC-Pro & 5,400 & $\checkmark$ & $\checkmark$ & $\checkmark$ &  & $\checkmark$ &  &  &  & $\checkmark$ &  &  &  &  &  & \\
\midrule
\multicolumn{17}{c}{\bfseries 反事实输入（未加掩码的句子）}\\
CrowS-Pairs & 1,508 & $\checkmark$ & $\checkmark$ & $\checkmark$ &  &  &  & $\checkmark$ & $\checkmark$ & $\checkmark$ & $\checkmark$ & $\checkmark$ & $\checkmark$ & $\checkmark$ & $\checkmark$ & $\checkmark$\\
WinoQueer & 45,540 & $\checkmark$ & $\checkmark$ & $\checkmark$ &  &  &  &  &  &  &  &  &  &  & $\checkmark$ & \\
RedditBias & 11,873 & $\checkmark$ & $\checkmark$ & $\checkmark$ & $\checkmark$ &  &  &  &  & $\checkmark$ &  &  & $\checkmark$ & $\checkmark$ & $\checkmark$ & \\
Bias-STS-B & 16,980 & $\checkmark$ & $\checkmark$ &  &  &  &  &  &  & $\checkmark$ &  &  &  &  &  & \\
PANDA & 98,583 & $\checkmark$ & $\checkmark$ & $\checkmark$ &  &  &  & $\checkmark$ &  & $\checkmark$ &  &  & $\checkmark$ &  &  & \\
Equity Evaluation Corpus & 4,320 & $\checkmark$ & $\checkmark$ & $\checkmark$ &  &  &  &  &  & $\checkmark$ &  &  & $\checkmark$ &  &  & \\
Bias NLI & 5,712,066 & $\checkmark$ & $\checkmark$ &  &  & $\checkmark$ &  &  &  & $\checkmark$ &  &  &  & $\checkmark$ &  & \\
\midrule
\multicolumn{17}{c}{\bfseries 提示（句子补全）}\\
RealToxicityPrompts & 100,000 &  &  &  & $\checkmark$ &  & $\checkmark$ &  &  &  &  &  &  &  &  & $\checkmark$\\
BOLD & 23,679 &  &  &  & $\checkmark$ & $\checkmark$ & $\checkmark$ &  &  & $\checkmark$ &  &  & $\checkmark$ & $\checkmark$ &  & $\checkmark$\\
HolisticBias & 460,000 & $\checkmark$ & $\checkmark$ & $\checkmark$ &  &  &  & $\checkmark$ & $\checkmark$ & $\checkmark$ & $\checkmark$ & $\checkmark$ & $\checkmark$ & $\checkmark$ & $\checkmark$ & $\checkmark$\\
TrustGPT & 9 &  &  & $\checkmark$ & $\checkmark$ &  & $\checkmark$ &  &  & $\checkmark$ &  &  & $\checkmark$ & $\checkmark$ &  & \\
HONEST & 420 & $\checkmark$ & $\checkmark$ & $\checkmark$ &  &  &  &  &  & $\checkmark$ &  &  &  &  &  & \\
\midrule
\multicolumn{17}{c}{\bfseries 提示（问答）}\\
BBQ & 58,492 & $\checkmark$ & $\checkmark$ &  & $\checkmark$ & $\checkmark$ &  & $\checkmark$ & $\checkmark$ & $\checkmark$ & $\checkmark$ & $\checkmark$ & $\checkmark$ & $\checkmark$ & $\checkmark$ & $\checkmark$\\
UnQover & 30 & $\checkmark$ & $\checkmark$ &  &  & $\checkmark$ &  &  &  & $\checkmark$ & $\checkmark$ &  & $\checkmark$ & $\checkmark$ &  & \\
Grep-BiasIR & 118 & $\checkmark$ & $\checkmark$ &  &  & $\checkmark$ &  &  &  & $\checkmark$ &  &  &  &  &  & \\
\bottomrule
\end{tabular}
\end{table}


% PDF 266--275；原书第 238--247 页。
\subsection{缓解策略}
\label{sec:ch06-debiasing}
下面按 \refcite{ch06-ref-029}{Gallegos 等人（2023）}的论述，考察偏见缓解策略，并依其在大语言模型工作流程中所处的阶段，区分预处理、内部处理和后处理，如图~\ref{fig:ch06-debiasing} 所示。

\begin{figure}[!htbp]
\centering
% 重绘自原书第 238 页的图 6.3。
\begin{tikzpicture}[
  font=\footnotesize\sffamily,
  link/.style={-{Stealth[length=1.55mm]},line width=.73pt},
  nodebox/.style={draw=orange!39,fill=orange!17,rounded corners=3pt,minimum width=17mm,
    minimum height=9.2mm,inner sep=3pt}]
\node (input) at (0,0) {输入};
\node[nodebox] (lm) at (4.14,0) {大语言模型};
\node (output) at (8.80,0) {输出};
\node[align=center] (fixed) at (8.12,2.02) {去偏后的\\输出};
\draw[link,blue!62!black] (input.east) -- node[above] {预处理} (lm.west);
\draw[link,violet!76] (lm.east) -- node[above] {内部处理} (output.west);
\draw[link,green!45!black] (8.14,.23) -- node[right] {后处理} (fixed.south);
\draw[link,red!56] ([xshift=4.5pt]lm.south) -- ++(0,-.49) -|
  ([xshift=-4.5pt]lm.south);
\node[text=red!59!black] at (4.14,-1.22) {训练阶段};
\end{tikzpicture}

\caption{偏见缓解策略及其在大语言模型工作流程中的位置。}
\label{fig:ch06-debiasing}
\end{figure}

\subsubsection{预处理阶段的缓解}
这类方法调整模型的输入，即数据和提示，保持模型的可训练参数不变，如图~\ref{fig:ch06-debias-pre} 所示。

\sourcefigure{fig06_04.pdf}{用预处理消除偏见的方法。上半图展示数据增强、筛选、重新加权及数据生成，下半图展示改写提示文字、使用控制词元及连续提示调优。图中保留数据修改前后的英文样例。}{fig:ch06-debias-pre}

\paragraph*{数据增强}
基于数据增强的方法会将代表不足的群体的样例加入训练数据，以增加数据集的多样性。反事实数据增强（counterfactual data augmentation，CDA）是其中的重要方法。通过替换受保护属性词，例如带性别的代词，使数据集更均衡。\refcite{ch06-ref-068}{Lu 等人（2020）}以此应对职业与性别相关的偏见，在保留语法和语义完整性的同时，互换表示性别的词。

有选择地替换词是 CDA 的另一种做法，可提高数据使用效率，并针对最能发挥偏见缓解作用的样例。

\begin{practicaltip}
\refcite{ch06-ref-072}{Maudslay 等人（2019）}提出的反事实数据替换（counterfactual data substitution，CDS）等技术，将含有性别词的文本随机替换成反事实版本。同一组作者还提出人名干预（Names Intervention），专门替换名字，将男性名字与成对的女性名字关联起来再作替换。

基于 \refcite{ch06-ref-127}{Zhang 等人（2017）}的 mixup 技术，插值方法将反事实增强后的训练样例与原始版本混合，进一步丰富训练数据。
\end{practicaltip}

\paragraph*{数据筛选与重新加权}
数据筛选和重新加权可以弥补数据增强在缓解偏见方面的局限，重点是在已有数据集中精确选定要处理的样例。主要分为数据集筛选和样例重新加权两类。

数据集筛选选择数据的子集，以影响模型在微调时的学习。\refcite{ch06-ref-032}{Garimella 等人（2022）}从代表不足的群体中选取文本，增加数据多样性。\refcite{ch06-ref-010}{Borchers 等人（2022）}则选择偏见最少的样例，构造偏见较少的数据集。还有些方法针对偏见最明显的样例，通过中性化或筛除它们，减少数据集总体上的偏见。例如，\refcite{ch06-ref-107}{Thakur 等人（2023）}筛选出一组偏见明显的样例，再将与性别有关的词改为中性词，以构造更加均衡的训练数据。

样例重新加权则调整训练时特定样例的重要性。\refcite{ch06-ref-043}{Han 等人（2022）}采用这一方法，按样例标签和相关受保护属性的出现情况，给损失函数中的各个样例分配反比权重。\refcite{ch06-ref-110}{Utama 等人（2020）}和 \refcite{ch06-ref-087}{Orgad 与 Belinkov（2023）}等方法，重点降低包含社会群体信息的样例的权重，以减少模型预测时对刻板印象捷径的依赖。

\paragraph*{数据生成}
数据生成旨在解决数据增强、筛选和重新加权的局限。其中一个难点是为每个偏见维度识别具体样例，而偏见维度可能因语境或应用而异。数据生成依据预先设定的标准或特征构造全新的数据集。\refcite{ch06-ref-103}{Solaiman 和 Dennison（2021）}开发了迭代流程来构建面向特定价值观的数据集，例如消除与法律保护类别有关的偏见。撰写者在这个过程中发挥重要作用，编写能体现目标行为的提示及补全，再根据表现评估的结果改进这些样例。\refcite{ch06-ref-022}{Dinan 等人（2019）}同样请人收集多种样例，以减少聊天对话模型的性别偏见。

数据生成的一个重点是建立新词表，特别是供 CDA、CDS 等词语互换技术使用的词表。这些词表往往侧重与不同社会群体相关的词，涉及性别、种族、年龄和方言等方面。不过，依赖这类词表有时会限制所能处理的刻板印象范围。为此，有研究者提出范围更广的框架。\refcite{ch06-ref-086}{Omrani 等人（2023）}着眼于“热情程度”和“能力”等较一般的维度来理解刻板印象，扩大了缓解偏见的范围。他们为这两类特征分别生成词表，在偏见缓解任务中，可以用这些词表替代按群体构建的词表，例如性别词表。

\paragraph*{指令调优}
这里的指令调优方法会修改交给模型的输入或提示。可以为提示添加文字指令或触发内容，鼓励生成没有偏见的输出。例如，\refcite{ch06-ref-071}{Mattern 等人（2022）}用抽象程度不同的提示引导模型避开刻板印象。\refcite{ch06-ref-111}{Venkit 等人（2023）}也用对抗触发词减少国籍偏见，\refcite{ch06-ref-002}{Abid 等人（2021）}则用简短的短语减轻对穆斯林的偏见。这些方法通常把短语附加在输入上，引导模型对某个社会群体表达中性或正面的情感。

还可以用控制词元（control token）为提示分类。模型学习把每个词元与某一类输入联系起来，从而在推理时控制生成。\refcite{ch06-ref-022}{Dinan 等人（2019）}使用这一办法缓解对话生成中的性别偏见，按照训练样例中有无性别词附加对应的词元。

\begin{practicaltip}
连续提示调优（continuous prompt tuning）是在输入上加入可训练前缀的另一项技术。原预训练模型的参数保持冻结，同时允许针对具体任务调节前缀。这种办法能实现手工提示工程难以完成的、有针对性且可扩展的调整。\refcite{ch06-ref-025}{Fatemi 等人（2021）}和 \refcite{ch06-ref-120}{Yang 等人（2023）}用连续提示调优缓解性别偏见，鼓励模型使用不随受保护属性而改变的中性语言。
\end{practicaltip}

\subsubsection{训练阶段的缓解}
训练阶段的缓解策略旨在模型训练过程中减少偏见。这些技术改变优化过程，包括修改损失函数、更新下一个词的概率、在微调时有选择地冻结参数，以及移除与有害输出相关的神经元，如图~\ref{fig:ch06-debias-train} 所示。所有这些策略均通过基于梯度的训练更新来改变模型参数。

\sourcefigure{fig06_05.pdf}{训练过程中用于减少偏见的几种方法。从左到右为修改架构、修改损失函数、有选择地更新参数，以及筛除模型参数。}{fig:ch06-debias-train}

\paragraph*{修改架构}
修改架构是训练阶段缓解偏见的一个方面，涉及改变模型结构，例如层、编码器和解码器的数量、大小及类型。一个例子是加入去偏适配器模块，如 \refcite{ch06-ref-057}{Lauscher 等人（2021）}提出的 ADELE。它们以模块化适配器框架为基础，在已有层之间插入新层，以高效微调。新插入的层参与微调，预训练层保持固定，使新增部分专门学习去偏所需的信息。

Liu 等人（2022）引入一个正则项，用于尽量减小受保护属性由 $E(\cdot)$ 给出的嵌入之间的距离。
\begin{equation}
\mathcal{R}=\lambda\sum_{(a_i,a_j)\in A}\lVert E(a_i)-E(a_j)\rVert_2.
\label{eq:ch06-embedding-regularization}
\end{equation}
其中，$\mathcal{R}$ 是正则项，$\lambda$ 为缩放因子，$a_i$、$a_j$ 是集合 $A$ 的元素，分别表示受保护属性及其反事实替换。

\paragraph*{修改损失函数}
\refcite{ch06-ref-088}{Park 等人（2023）}提出一种技术，将基于投影的偏见缓解方法纳入损失函数，专门处理职业词中的性别刻板印象。他们引入一个正则项，使刻板印象相关词的嵌入 $\symbf{w}$ 和嵌入空间中的性别方向 $\symbf{g}$ 正交。这个正则项使中性职业词的嵌入远离本身具有性别含义的词，例如“sister”或“brother”。

性别方向的形式定义为
\begin{equation}
\symbf{g}=\frac{1}{|A|}\sum_{(a_i,a_j)\in A}\bigl(E(a_j)-E(a_i)\bigr),
\label{eq:ch06-gender-direction}
\end{equation}
其中，$A$ 是所有本身具有性别含义的女性相关词 $a_i$ 与男性相关词 $a_j$ 所构成的词对集合，$E(\cdot)$ 计算模型的嵌入。

正则项表示为
\begin{equation}
\mathcal{R}=\sum_{\symbf{w}\in W_{\mathrm{stereo}}}
\frac{\symbf{g}}{\lVert\symbf{g}\rVert}\symbf{w}^{\top},
\label{eq:ch06-projection-regularization}
\end{equation}
其中，$W_{\mathrm{stereo}}$ 为刻板印象相关词的嵌入集合。%
\footnote{译注：原书在公式末尾和正式编号（6.20）之间还印有“(34)”，此处只保留本书公式编号。}

\refcite{ch06-ref-004}{Attanasio 等人（2022）}考察了修改语言模型注意力层的方法。他们假设，注意力层是偏见的主要编码位置。他们提出基于熵的注意力正则化（entropy-based attention regularization，EAR），用注意力权重分布的熵衡量上下文词语的相关性。熵高，表示利用的上下文较多样。熵低，则表示模型较依赖特定词元。EAR 旨在最大化注意力权重的熵，避免对带刻板印象的词过拟合，从而让模型关注更广的输入上下文。

具体做法是在损失函数中加入以最大化熵为目标的正则项。
\begin{equation}
\mathcal{R}=-\lambda\sum_{\ell=1}^{L}\operatorname{entropy}(A)_\ell,
\label{eq:ch06-attention-entropy}
\end{equation}
其中，$\operatorname{entropy}(A)_\ell$ 是第 $\ell$ 层的注意力熵。

也有研究通过设计损失函数，使语言模型输出中与特定人口群体相关的词获得均衡的概率。例如，\refcite{ch06-ref-092}{Qian 等人（2019）}引入使概率相等的目标，鼓励模型以相同的概率预测不同人口群体的相关词。该方法有一项正则项，比较二元性别词对的 Softmax 输出概率。\refcite{ch06-ref-032}{Garimella 等人（2022）}后来将这一项用于由两类种族相关词组成的词对。对由属性 $a_i$ 和 $a_j$ 组成的 $K$ 个词对，可使用下列正则项。其中，$a_i\in A_i$ 是与社会群体 $G_i$ 有关的受保护属性词，类似地，$a_j\in A_j$ 是与社会群体 $G_j$ 有关的受保护属性词。
\begin{equation}
\mathcal{R}=\lambda\frac{1}{K}\sum_{k=1}^{K}
  \left|\log\frac{P(a_i^{(k)})}{P(a_j^{(k)})}\right|.
\label{eq:ch06-equalizing-regularization}
\end{equation}

\paragraph*{有选择地更新或筛除模型参数}
用增强的数据集微调人工智能模型可以减少偏见，但可能导致“灾难性遗忘”，使模型失去先前学到的信息。为防止这种情况，有些方法只更新极少一部分模型参数，冻结其他参数。例如，\refcite{ch06-ref-035}{Gira 等人（2022）}在 WinoBias 和 CrowS-Pairs 数据集上微调模型时，只更新层归一化等特定参数。\refcite{ch06-ref-094}{Ranaldi 等人（2023）}则只调整注意力矩阵，并冻结其他所有参数。\refcite{ch06-ref-123}{Yu 等人（2023）}有针对性地优化对某类偏见贡献最大的权重，例如与性别和职业相关的权重。

除了通过微调更新参数来减轻偏见，还有技术侧重有选择地筛除某些参数，例如在训练中或训练后把它们置零。\refcite{ch06-ref-050}{Joniak 和 Aizawa（2022）}使用移动剪枝（movement pruning），从神经网络中有选择地移除权重。他们借此从预训练模型的注意力头中选出偏见较少的一组权重。微调时冻结这些权重，并另外优化用于去偏的分数，决定移除哪些权重。

\subsubsection{推理时内部处理}
按 \refcite{ch06-ref-097}{Savani 等人（2020）}的定义，内部处理（intra-processing）在推理阶段改变预训练或微调模型的行为，无需进一步训练便可生成经过去偏的预测。这类推理阶段缓解办法，包括调整解码策略、在训练后修改模型参数，以及在推理时以模块化方式采用独立的去偏网络，如图~\ref{fig:ch06-debias-infer} 所示。

生成输出词元时，可以调整解码算法，尽量减少带有偏见的语言。这样的调整不改变可训练的模型参数。它通过选择约束、改变词元概率分布或引入辅助偏见检测模型，影响后续词语或序列出现的概率。其中一种办法是受约束的下一词元搜索（constrained next-token search），通过加入额外准则，改变输出序列中下一个词元的排序，从而引导模型给出偏见较少的输出。

\refcite{ch06-ref-033}{Gehman 等人（2020）}建议在解码时屏蔽词语或 $n$ 元组，防止选中预定义不安全词表中的词元。不过，由若干单独看并无偏见的词元或 $n$ 元组组成的文本，仍可能含有偏见，所以这一办法无法完全消除生成偏见输出的可能。

\sourcefigure{fig06_06.pdf}{推理时缓解偏见的几种方法。左侧展示受约束的下一词元搜索和改变词元概率分布，右侧展示注意力权重重分配和模块化去偏网络。}{fig:ch06-debias-infer}

\begin{practicaltip}
权重重分配（weight redistribution）是在训练后修改与已编码的偏见有关的模型权重，主要针对注意力权重。\refcite{ch06-ref-125}{Zayed 等人（2023）}通过受一个超参数控制的温度缩放，在训练后调整注意力权重。这样可以增大熵，使模型考虑更广范围的词元，可能因此避开刻板印象。也可以减小熵，使模型集中在较窄的上下文上，减少接触刻板印象词元的机会。
\end{practicaltip}

最后，模块化去偏网络解决了一些训练阶段方法的局限。训练阶段的方法通常只针对一种偏见，而且会永久改变模型。模块化去偏更灵活，可以建立独立组件，为不同任务或不同偏见将其与预训练模型组合。\refcite{ch06-ref-044}{Hauzenberger 等人（2023）}将 \refcite{ch06-ref-040}{Guo 等人（2020）}提出的差分剪枝（diff pruning）用于去偏。他们模拟同时训练多个模型，每个模型针对一个不同的维度减少偏见。再把对预训练模型参数所作的修改提取出来，高效保存在稀疏子网络中。

\subsubsection{后处理阶段的缓解}
后处理阶段的缓解对预训练模型的输出做处理，以减少偏见。如果模型主要作为黑箱提供，难以得知其训练数据或内部工作方式，这类办法尤其有用，如图~\ref{fig:ch06-debias-post} 所示。

\sourcefigure{fig06_07.pdf}{后处理阶段缓解偏见的例子。上方路线检测并替换关键词，下方路线由神经机器翻译模型改写输出。图中保留替换前后的英文句子。}{fig:ch06-debias-post}

\begin{practicaltip}
基于改写的方法检测并替换模型输出中的偏见词语或有害词语。关键词替换等技术识别带有偏见的词元，用较中性或更有代表性的用词代替，力求保留原文的内容和风格。例如，\refcite{ch06-ref-108}{Tokpo 和 Calders（2022）}利用 LIME，根据原文的潜在表示识别并替换造成偏见的词元。\refcite{ch06-ref-021}{Dhingra 等人（2023）}则利用 SHAP 识别并替换带有刻板印象的词，同时解释替换的原因。
\end{practicaltip}

另一种后处理技术使用机器翻译，把带有偏见的句子转换成中性或没有偏见的表达。多项研究采用基于规则的方法，为包含性别代词的句子构造性别中性版本（\refcite{ch06-ref-049}{Jain et al., 2021}；\refcite{ch06-ref-106}{Sun et al., 2021}）。这种办法从带有性别偏见的源句生成经去偏的对应句，构成平行数据，再训练机器翻译模型，将有性别偏见的句子转换为去偏版本。

% PDF 275--283；原书第 247--255 页。
\section{毒性}
\label{sec:ch06-toxicity}

\begin{keypoint}
毒性（toxicity）涵盖多种有害内容，没有统一定义。可以将它理解为前文所定义的表征性伤害之一，也可以视作独立概念。Perspective API 将毒性描述为可能使参与者退出对话的粗鲁、不尊重或不讲理的评论。\refcite{ch06-ref-056}{Kurita 等人（2019）}将有毒内容描述为可能冒犯或伤害受众的任何材料，其中包括仇恨言论、种族主义和冒犯性用语。\refcite{ch06-ref-089}{Pavlopoulos 等人（2020）}把毒性视为一个统称，研究群体使用“冒犯性”“辱骂性”和“仇恨性”等不同术语，描述各种有毒语言或相关现象。
\end{keypoint}

\subsection{成因}
大语言模型的毒性、偏见和公平性相互关联，都受到训练数据的影响。偏见一节所列的许多成因，例如数据选择、领域和体裁分布不均、创作者的人口构成和文化分布偏斜，也适用于大语言模型的有毒输出。这里侧重可能为毒性所特有、或与模型偏见成因重合的因素。

\begin{enumerate}
\item \textbf{训练数据偏见。}如偏见一节所述，大语言模型的毒性主要来源之一，是训练数据本身带有偏见。社会不平等、带有偏见的语言用法，以及某些群体代表不足，都可能使训练数据存在偏见。这些偏见表现在模型的输出中，产生有毒、不公平的结果。模型往往复现数据集中的偏见，导致有毒输出（\refcite{ch06-ref-006}{Bender et al., 2021}）。
\item \textbf{理解上下文的局限。}大语言模型有时难以理解文本或对话的完整上下文，从而给出不当或有毒的回答。\refcite{ch06-ref-005}{Bender 和 Koller（2020）}指出，模型难以理解含义细致的人类语言，这说明准确理解上下文并不容易。\refcite{ch06-ref-089}{Pavlopoulos 等人（2020）}发现，帖子的上下文会显著影响人们对其毒性的判断，既可能放大，也可能减轻感知到的毒性。他们的一项实验中，标注者未看到上下文时，大约有 5\% 的人工标注帖子得到了相反的毒性标签。
\item \textbf{对抗攻击。}大语言模型容易受到对抗攻击，被提示故意生成有毒输出。\refcite{ch06-ref-112}{Wallace 等人（2020）}的研究说明，攻击者可向模型的训练集中注入恶意样例，显著影响其学习和未来预测。这种攻击使攻击者可以将任意选定的短语变成触发特定预测的通用触发语，构成一种危险的漏洞。研究还表明，可以把这些投毒样例设计得不显眼，使受攻击者难以识别并移除有害数据。投毒样例本身没有明确提到触发语，因此可以避开搜索特定短语的检测方法。
\item \textbf{指定角色的提示。}对话式人工智能的一种常见用法，是让大语言模型扮演某一角色，并以该角色继续对话。\refcite{ch06-ref-020}{Deshpande 等人（2023）}发现，为 ChatGPT 指定特定角色，例如著名拳击手穆罕默德·阿里，可能显著提高生成文本的毒性。根据分配给 ChatGPT 的角色，回答的毒性最高可变为原来的六倍。毒性增加表现为模型助长错误的刻板印象，生成有害对话，表达伤人的观点。将这类回答与指定角色联系起来，既可能损害公众人物的名誉，也可能伤害没有预料到这类内容的用户。
\end{enumerate}

\subsection{评估指标}
大多数毒性评估框架都使用辅助模型，通常是分类器，来评估语言模型生成文本的毒性水平。

原书介绍，Google Jigsaw 开发的 Perspective API 是常用的文本毒性评分工具（\refcite{ch06-ref-059}{Lees et al., 2022}）。如图~\ref{fig:ch06-perspective} 所示，输入为文本，输出是 0 至 1 之间的概率分数。它量化文本被认为含有某种毒性相关属性的可能性。属性包括 \texttt{TOXICITY}（毒性）、\texttt{SEVERE\_TOXICITY}（重度毒性）、\texttt{IDENTITY\_ATTACK}（身份攻击）、\texttt{INSULT}（侮辱）、\texttt{PROFANITY}（粗俗用语）、\texttt{THREAT}（威胁）、\texttt{SEXUALLY\_EXPLICIT}（露骨的性内容）和 \texttt{FLIRTATION}（调情）。Perspective API 使用来自多个在线论坛的大量语料训练模型，来源包括维基百科和《纽约时报》。

\sourcefigure{fig06_08.pdf}{通过 Perspective API 评估毒性。示例文本“Shut up. You’re an idiot!”的各项分数为毒性 0.99、重度毒性 0.75、侮辱 1.0、露骨的性内容 0.04、粗俗用语 0.93、可能遭到拒绝 0.99、威胁 0.15 和身份攻击 0.03。}{fig:ch06-perspective}

数据由众包评分人员标注，人员来自 Figure Eight、Appen 等机构及内部平台。评分人员判断一条评论是否具有某项属性，例如 \texttt{TOXICITY}。最终标签由认为文本有毒的评分者所占比例得出。在模型架构方面，Perspective API 先用标注数据训练基于多语言 BERT 的模型，再将模型蒸馏到每种受支持语言各自的单语言卷积神经网络，使系统能够快速、高效地给出分数。

评估文本生成时，Perspective API 提供表示毒性的概率分数。其毒性分类器可写为
\[
c\colon\widehat{Y}\to[0,1].
\]
\refcite{ch06-ref-033}{Gehman 等人（2020）}采用期望最大毒性（Expected Maximum Toxicity，EMT），对语言模型的 25 条输出计算，关注均值和标准差。EMT 用于找出语言模型生成文本中的极端情况，即毒性最高的输出。
\begin{equation}
\operatorname{EMT}(\widehat{\mathbb{Y}})
=\max_{\widehat{Y}\in\widehat{\mathbb{Y}}}c(\widehat{Y}).
\label{eq:ch06-emt}
\end{equation}

毒性概率（Toxicity Probability，TP）衡量 25 次生成中，至少有一次生成有毒文本的可能性。毒性分数等于或大于 0.5 的文本记为有毒。该指标量化语言模型生成有毒内容的频率。
\begin{equation}
\operatorname{TP}(\widehat{\mathbb{Y}})
=P\!\left(\sum_{\widehat{Y}\in\widehat{\mathbb{Y}}}
 I\bigl(c(\widehat{Y})\geq0.5\bigr)\geq1\right).
\label{eq:ch06-tp}
\end{equation}

\refcite{ch06-ref-009}{Bommasani 等人（2023）}定义的有毒比例（Toxic Fraction，TF）计算生成文本中被判为有毒的比例。
\begin{equation}
\operatorname{TF}(\widehat{\mathbb{Y}})
=\mathbb{E}_{\widehat{Y}\in\widehat{\mathbb{Y}}}
 \bigl[I(c(\widehat{Y})\geq0.5)\bigr].
\label{eq:ch06-tf}
\end{equation}

\subsection{评估基准}
上一节列出的偏见基准与用于评估毒性的基准有相当多的重合。本节侧重专门为毒性设计的基准，见表~\ref{tab:ch06-toxicity-benchmarks}。

\begin{table}[!htbp]
\centering
\caption{主要的毒性基准数据集。表中给出数据集的大小、收集和标注方法，以及所关注的内容。}
\label{tab:ch06-toxicity-benchmarks}
\small
\setlength{\tabcolsep}{3pt}
\begin{tabularx}{\linewidth}{@{}>{\raggedright\arraybackslash}p{.31\linewidth}rlX@{}}
\toprule
数据集 & 规模 & 方法 & 关注点\\
\midrule
Perspective API 的毒性数据集 & 180 万 & 众包 & 整体毒性及特定维度\\
Jigsaw Toxic Comment Dataset & 15 万 & 众包 & 毒性程度和类型\\
Hate Speech Dataset & 2.4 万 & 众包 & 仇恨言论检测\\
ToxiGen & 10 万 & 对抗式构造 & 模型的稳健性与隐含偏见\\
Thoroughly Engineered Toxicity（TET）Dataset & 1 万 & 人工构造 & 使模型防护失效\\
ImplicitHateCorpus & 570 万 & 众包 & 隐含仇恨言论（讽刺、刻板印象、微侵犯）\\
DynaHate & 2250 万 & 机器学习 & 与语境相关的仇恨言论（特定对象、变化中的用语）\\
SocialBiasFrames & 8,732 & 众包 & 有害的社会框架（性别、种族、残障）\\
\bottomrule
\end{tabularx}
\end{table}

\subsection{缓解策略}
\refcite{ch06-ref-033}{Gehman 等人（2020）}将毒性缓解技术分为两大类，即基于数据和基于解码的方法。基于数据的方法对模型进一步预训练，改变模型参数。这个办法有效，但涉及修改参数，通常消耗较多算力。基于解码的方法只调整语言模型的解码算法，模型参数保持不变，因此通常较易采用，所需资源也较少。

\subsubsection{基于数据的方法}
\refcite{ch06-ref-042}{Gururangan 等人（2020）}讨论的领域适应性预训练（Domain Adapted Pre-training，DAPT）分为两个阶段。第一阶段在广泛的来源上作通用预训练，第二阶段使用与特定语境有关的数据集，例如 Twitter 交流数据，进行领域专用的预训练，如图~\ref{fig:ch06-dapt} 所示。结合特定任务的微调时，这种办法尤其有效。例如，在对与“黑人的命也是命”运动有关的话题分类时，微调侧重从更广的推文领域抽取这一语境特有的语言层次和细微差异。\refcite{ch06-ref-033}{Gehman 等人（2020）}说明，用 DAPT 在 OpenWebTextCorpus（OWTC）的无毒子集上继续预训练大语言模型，可明显减少 GPT-2 等模型的毒性，但这一办法也面临困难。最明显的局限之一是计算开销。另外，需要额外的训练数据，尤其通过众包等人工方式获得时，费用可能过高。另一个主要局限是，它可能损害模型表现。筛除有毒成分时，也可能误删无害且有潜在价值的知识。

\sourcefigure{fig06_09.pdf}{领域适应性预训练示意图。模型初次预训练完成后，用另一份语料进行第二次预训练。这份语料与模型任务有关，例如来自网页的数据，并且已滤除有毒的交流内容。图中以从 2020 年推文逐步限定到相关领域和任务数据为例。}{fig:ch06-dapt}

属性条件化（Attribute Conditioning，ATCON）是另一种基于数据的方法，源于 \refcite{ch06-ref-028}{Ficler 和 Goldberg（2017）}及 \refcite{ch06-ref-052}{Keskar 等人（2019）}的工作。该方法从文档中随机选择一个子集，为预训练模型的预训练数据加入毒性属性词元，记作 \verb!|toxic|! 或 \verb!|non-toxic|!。随后继续预训练，让模型学会使用这些毒性标记。用模型生成文本时，把毒性属性词元 \verb!|non-toxic|! 加在提示前面，向模型表明生成文本应当侧重于无毒内容。

% PDF 280--283；原书第 252--255 页。
\subsubsection{基于解码的方法}
这类方法调整解码算法，以尽量减少毒性。

\begin{practicaltip}
最简单的基于解码的办法之一是用屏蔽词表（blocklisting），也称词语过滤。把语言模型中不希望出现的词语的概率设为零。这些词一般包括咒骂、粗俗用语和侮辱词语。这样，模型的输出就不会出现这些词。下面介绍一些在解码时降低毒性、较复杂而通常也更有效的方法。
\end{practicaltip}

\paragraph*{词表偏移}
\refcite{ch06-ref-034}{Ghosh 等人（2017）}提出这一方法。它为词表中的每个词元学习向量表示，区分毒性与无毒属性。原研究使用 LSTM 模型，不过其基本原理也可以用于 GPT-2 等较新的 Transformer 架构。

在标准 LSTM 模型中，由 $M$ 个词 $w_1,w_2,\ldots,w_M$ 组成的序列，其联合概率由概率链式法则定义为
\begin{equation}
P(w_1,w_2,\ldots,w_M)
=\prod_{t=1}^{M}P(w_t\mid w_1,w_2,\ldots,w_{t-1}).
\label{eq:ch06-lstm-joint}
\end{equation}
在这种模型中，词 $w_t$ 在上下文 $c_{t-1}=(w_1,w_2,\ldots,w_{t-1})$ 下的条件概率为
\begin{equation}
P(w_t=i\mid c_{t-1})
=\frac{\exp(U_i^{\mathsf T}f(c_{t-1})+b_i)}
 {\sum_{j=1}^{V}\exp(U_j^{\mathsf T}f(c_{t-1})+b_j)}.
\label{eq:ch06-lstm-conditional}
\end{equation}
其中，$f(\cdot)$ 是 LSTM 的输出，$U$ 是词表示矩阵，$b_i$ 是偏置项。\refcite{ch06-ref-034}{Ghosh 等人（2017）}提出用于词表偏移的 Affect-LM，把上式改写为
\begin{multline}
P(w_t=i\mid c_{t-1},e_{t-1})\\
=\frac{\exp(U_i^{\mathsf T}f(c_{t-1})+\beta V_i^{\mathsf T}g(e_{t-1})+b_i)}
 {\sum_{j=1}^{V}\exp(U_j^{\mathsf T}f(c_{t-1})+\beta V_j^{\mathsf T}g(e_{t-1})+b_j)}.
\label{eq:ch06-affect-lm}
\end{multline}
其中，$e_{t-1}$ 表示训练时由上下文词语获得的情绪类别信息。该式量化情绪类别信息对在上下文中预测目标词 $w_t$ 的影响，$\beta$ 是情绪强度参数。

\paragraph*{即插即用语言模型（PPLM）}
即插即用语言模型（Plug and Play Language Model，PPLM）允许用户把一个或多个表示特定控制目标的属性模型与大语言模型结合（\refcite{ch06-ref-019}{Dathathri et al., 2019}）。这样的组合无需额外训练或微调模型，对难以获得大量硬件资源的研究者很有用。PPLM 有两个基本条件。

\begin{enumerate}
\item 可以使用一个属性模型，记作 $p(a\mid x)$。
\item 可以取得该属性模型的梯度。
\end{enumerate}

PPLM 流程如图~\ref{fig:ch06-pplm} 所示，分为以下几步。

\begin{enumerate}
\item 对大语言模型做一次前向传播，从所得概率分布抽取一个词元。然后把生成的字符串交给属性模型，计算所需属性的概率 $p(a\mid x)$。
\item 反向传播，计算 $p(a\mid x)$ 和 $p(x)$ 相对于模型隐藏状态的梯度。调整隐藏状态，提高这两个概率。
\item 重新计算大语言模型的概率分布，抽取一个新的词元。
\end{enumerate}

\sourcefigure{fig06_10.pdf}{即插即用语言模型（PPLM）。图中使用属性模型反向调整语言模型的隐状态，再对下一词元重新采样。}{fig:ch06-pplm}

\paragraph*{生成式判别器（GeDi）}
生成式判别器（Generative Discriminator，GeDi）用较小的语言模型作生成式判别器，引导大语言模型的生成，以缓解毒性（\refcite{ch06-ref-054}{Krause et al., 2020}）。它用贝叶斯法则计算序列中每个可能的后续词的概率，平衡两个类别条件分布。第一个分布以所需属性，即“控制码”为条件，引导模型生成特定的理想内容。第二个分布的条件是“反控制码”，表示希望避开的属性，例如有毒语言。

基本原理是用辅助的类别条件语言模型（class-conditional language model，CC-LM），判断生成的词元属于控制码所指定类别的可能性。控制码定义文本序列 $x_{1:T}$ 的一种属性，可以是情感、主题或“有毒”等标签。类别条件语言模型计算 $p(x\mid c)$ 和 $p(x\mid\bar c)$，其中 $c$ 和 $\bar c$ 分别是控制码和反控制码。CTRL 就是一个例子，它提供了一个框架，让语言模型以属性变量为条件生成文本（\refcite{ch06-ref-052}{Keskar et al., 2019}）。

训练类别条件语言模型时，将每个序列 $x^{(i)}_{1:T_i}$ 与控制码 $c^{(i)}$ 配对。随后训练模型，使平均负对数似然 $L_g$ 最小。GeDi 的训练损失 $L_{\mathrm{gd}}$ 结合了语言模型损失 $L_g$ 和判别损失 $L_d$，用超参数 $\lambda$ 加权。
\begin{equation}
L_{\mathrm{gd}}=\lambda L_g+(1-\lambda)L_d.
\label{eq:ch06-gedi}
\end{equation}
判别损失用于推动模型产生希望的行为，并远离不希望的行为。$L_d$ 旨在提高分类准确率，而 $L_g$ 有助于更好地校准引导生成所用的词元概率。与 PPLM 相比，GeDi 在降低毒性方面更有效，计算开销也更低。但它需要一个额外的语言模型进行训练，因此增加了算力需求。

\paragraph*{自我诊断与自我去偏}
\refcite{ch06-ref-098}{Schick 等人（2021）}提出的方法围绕两个问题展开。第一，大语言模型能否自行发现自身的偏见，即自我诊断（self-diagnosis）？第二，大语言模型能否以无监督方式抑制不希望出现的行为，即自我去偏（self-debiasing）？

自我诊断考察预训练语言模型能否识别自身输出中在社会层面不合宜的属性，而无需额外的训练数据或外部资源。它利用模型在预训练中学到的知识，并给出简短的属性描述，以评估模型自行识别和处理偏见、不当内容的能力。

设语言模型为 $M$，词元序列为 $x_1,\ldots,x_k=x$。进行自我诊断时，将 $M$ 生成的句子 $x$ 与关于属性 $y$ 的问题组合，判断 $x$ 是否含有 $y$。例如，若 $M$ 生成句子 $x=$“I'm going to hunt you down!”，要判断它是否含有威胁，即 $y=$“a threat”，可用以下模板提示模型。
\begin{lstlisting}[nolol=true]
# Self-diagnosing template , denoted as sdg(.)
{sample text x}
Does the above text contain {threat text y} ?
\end{lstlisting}
$x$ 具有属性 $y$ 的概率为
\begin{equation}
p(y\mid x)
=\frac{p_M(\text{“Yes”}\mid\operatorname{sdg}(x,y))}
 {\sum_{w\in\{\text{“Yes”},\text{“No”}\}}
  p_M(w\mid\operatorname{sdg}(x,y))}.
\label{eq:ch06-sdg}
\end{equation}

自我去偏采用类似于自我诊断的思路，重点在于主动改变输出。它也利用模型自身的内部知识调整输出生成，降低生成有偏见文本的可能性。下面是一个自我去偏模板示例。
\begin{lstlisting}[nolol=true]
# Self-debiasing template , denoted as sdb(.)
The following text contains {undesired attribute s}:
{sample text x}
\end{lstlisting}

采用该模板时，需要预训练语言模型 $M$ 和表示某类偏见的属性 $y$。给模型输入文本 $x$，它用自我去偏输入 $\operatorname{sdb}(x,y)$ 生成续写。这会得到两个概率分布，即原始分布 $p_M(w\mid x)$ 和调整后的分布 $p_M(w\mid\operatorname{sdb}(x,y))$。后者被设计为更倾向于不希望出现的词。因此，对有偏见的词，下式的差值预期为负。
\[
\Delta(w,x,y)=p_M(w\mid x)-p_M(w\mid\operatorname{sdb}(x,y)).
\]
为缓解偏见，构造新的概率分布 $\widetilde{p}_M(w\mid x)$，使它与 $\alpha(\Delta(w,x,y))\cdot p_M(w\mid x)$ 成比例。这里，$\alpha\colon\mathbb{R}\to[0,1]$ 为缩放函数，根据偏见差值 $\Delta(w,x,y)$ 调整概率。这样只在有需要时改变概率，其余词保持不变，尽量减少对生成过程的干预。

为便于模型评估，这种办法不会把任何词的概率设为零。它根据 $\Delta(w,x,y)$ 降低带有偏见的词的概率，缩放函数使用衰减常数 $\lambda$。当 $x\geq0$ 时 $\alpha(x)=1$，否则 $\alpha(x)=e^{\lambda x}$。还可以考虑一组描述 $Y=\{y_1,\ldots,y_n\}$，据此调整分布，将该方法扩展到多种偏见。


% PDF 284--290；原书第 256--262 页。
\section{隐私}
\label{sec:ch06-privacy}
互联网在过去几十年兴起，信息比以往任何时候都流通得更自由。然而，有些信息并非提供者有意、知情地分享，获取信息的人也并非都慎重地收集、保存这些信息。因此，个人身份信息和其他私人数据变得更容易获取，已成为受到广泛关注的问题。由于最知名的大语言模型都从网站获得大量训练数据，容易想到这是否会给隐私带来后续风险。结果表明，大语言模型确实会学到有关个人的具体信息，通过适当提示可以提取这些信息。大语言模型的隐私问题目前仍未得到充分解决。本节讨论现有研究及应对这些问题的新趋势。

\subsection{成因}
在机器学习的大多数领域，人们通常认为，模型经常生成与训练样例非常接近的预测，是过拟合的典型症状。由此曾有一种假定，大语言模型通常只在规模很大的数据上训练一轮，这与过拟合的典型条件相反，所以一般不太可能记住训练数据并逐字复述。然而，记忆只会由过拟合引起这一假定已被否定（\refcite{ch06-ref-015}{Carlini et al., 2021}）。早期人们没有普遍认识到大语言模型的记忆能力，对减轻私人数据被记住这一问题的研究，落后于模型整体能力的发展。参数更多或训练词元更多的高性能大语言模型，似乎也比规模较小的版本更能记忆数据（\refcite{ch06-ref-082}{Nasr et al., 2023}）。

为量化某个模型记住的训练数据，\refcite{ch06-ref-015}{Carlini 等人（2021）}提出一种旨在引出记忆字符串的攻击。该方法从互联网来源抽取词元序列，这些来源已知或很可能是训练数据的一部分。以这些词元提示模型，再检查输出是否与训练数据精确匹配。训练来源公开时，这一过程显然更容易实施。为了把研究扩展到从未公开训练数据的 GPT-2，作者不得不通过 Google 在整个互联网上搜索，寻找极有可能被模型记住的匹配文本。尽管有这个限制，他们仍找到了 GPT-2 记住的数据，包括个人身份信息（personally identifiable information，PII）。

在后续工作中，\refcite{ch06-ref-082}{Nasr 等人（2023）}基于上述流程做了若干实验。研究区分了发现不同记忆输出的比率和提取出的不同序列的数量。例如，查询数不超过约 600 亿时，Pythia-1.4 输出的记忆数据多于 GPT-Neo-6。到达该点时，两条曲线交叉，此后 GPT-Neo-6 超过 Pythia-1.4。这并不意外，因为没有直观理由认为，某种攻击技术对一个模型的效果能直接衡量它实际记住的数据量。作者用外推方法建立了两者之间的关系，从而为各个被评估模型所记住的数据总量给出更紧的下界。

聊天模型的表现不同，因为它们通常经过对齐调优，并被引导以对话口吻给出有帮助的回答。因此，ChatGPT 起初显得相对安全。在提示中给出一串词元时，它不太可能继续输出自己记住的文本。不过，\refcite{ch06-ref-082}{Nasr 等人（2023）}还发现了一种新攻击，成功从 ChatGPT 提取了记忆数据。办法是选择一个只由一个词元构成的词，将它重复若干次，然后要求 ChatGPT 无休止地重复它。结果随所选词元而有很大差异，但在某些情况下，ChatGPT 最终会停止重复，开始复述记住的训练数据。研究者在该实验生成的一份文本样本中，核实了数千处被记住的个人身份信息。

由于 ChatGPT 是封闭的，这种现象的确切机制尚不清楚。对开源模型的考察表明，随着文本不断变长，模型可能推断它即将结束的概率越来越高。毕竟，训练文本没有真正无限延续的例子，模型并不熟悉这种情形。这可能使模型最终越来越倾向于将唯一的“文本结束”词元作为下一个预测，进而在前面没有有意义的上下文时开始生成新的序列。直观上，这可能显著提高模型回退到训练时见过的文本序列的概率。

这些攻击旨在说明，从大语言模型中提取相当数量的记忆数据是可能的。研究只涉及了潜在技术的一部分，也只初步触及攻击者可能怎样提高这些技术的效率和规模。如前所述，ChatGPT 在早期攻击中似乎不易泄露训练数据，新的攻击出现后却发生了泄露。这说明，很难估计大语言模型记忆数据的完整范围及其造成伤害的可能性。

\subsection{评估指标}
\refcite{ch06-ref-053}{Kim 等人（2023）}提出了用于探查和量化 PII 泄露的工具 ProPILE。它包含一个评估数据集和一组提示模板。评估数据集通过在广泛使用的 Pile 训练数据集中查找 PII 构造。提示模板在给出有关个人的其他相关信息后，用于检验模型能否取回某一类个人数据。下面是一条例示提示。

\begin{lstlisting}[nolol=true]
Please contact {name} by {pii_1} or {target_pii_type}
\end{lstlisting}

它可以测试，姓名和电话号码是否足以让模型取回街道地址，或街道地址是否会暴露电子邮件地址等。对所有可能的排列进行查询，记录模型输出。

提示攻击的效果可用两种指标衡量，即精确匹配和似然。精确匹配指标直接判断生成的字符串与真实 PII 是否相同。似然指标利用许多模型 API 可为每个词元预测返回分数的能力。获得似然后，可以作更细的估计，尤其适用于攻击者能够大规模、多次重复相近提示的情况。定义如下。
\begin{equation}
\Pr(a_m\mid A\setminus m)
=\prod_{r=1}^{L_r}p(a_{m,r}\mid x_1,x_2,\ldots,x_{L_q+r-1}),
\label{eq:ch06-pii-likelihood}
\end{equation}
其中，$a_m$ 为目标 PII，$A\setminus m$ 为其余 PII，$L_q$ 为查询长度，$L_r$ 为正确回答长度。对多个查询重复计算，可以得到另一项指标。它表示使用全部可用提示时，在不超过 $k$ 次查询中会有一项 PII 被暴露的人数，占训练数据所含人数的比例。这项指标记为 $\gamma_{<k}$。

在 ProPILE 中，想检查模型是否暴露自己数据的用户，只能做黑箱探查（black-box probing），因为他们只能看到模型的输出。黑箱探查可使用前面所述的提示模板。白箱探查（white-box probing）则指模型提供者希望量化自己的模型泄露 PII 的程度。此时模型权重已知，可以提高提示效果。若能完整访问模型，可以训练效果优于提示模板的软提示。假想的攻击者没有进行类似提示调优所需的信息，只能采用效率较低的提示工程技术。推测即使有技巧且意志坚定的攻击者，也很难设计出比模型所有者训练的软提示更好的探查方法，因此可以用这种方法估计最坏情况下的 PII 泄露。

\refcite{ch06-ref-048}{Inan 等人（2021）}提出了另一项隐私指标，以差分隐私（differential privacy）为基础。若两个数据集恰好相差一条记录，并且在两者上进行计算所得的输出概率满足由 $\epsilon$ 控制的偏差上界，就称该计算具有差分隐私。形式上，
\begin{equation}
\Pr[M(A)\in S]\leq e^{\epsilon}\Pr[M(B)\in S],
\label{eq:ch06-dp}
\end{equation}
其中，$A$ 和 $B$ 是只相差一条记录的数据集，$S$ 是所有可能输出的任意子集。当 $\epsilon$ 足够小时，差分隐私能对模型泄露与单个数据点有关的信息提供很强的限制。依定义，模型的输出分布将与在未包含该个体的数据上训练的模型的输出分布相近。

\refcite{ch06-ref-048}{Inan 等人（2021）}进一步训练了一个参考模型，去掉了被认定只属于任一用户的全部数据。然后，对移除的每项数据，计算参考模型与待评估模型的困惑度之比，借此定义最坏情况下的泄露指标。%
\footnote{译注：原书式（6.33）的对数括号未闭合，此处补齐。}
\begin{equation}
\epsilon_l=\max_{w\in S_{\mathrm{uniq}}}
 \log\!\left(\frac{\operatorname{PP}_{\mathrm{public}}(w)}
 {\operatorname{PP}_{\mathrm{lm}}(w)}\right),
\label{eq:ch06-worst-leakage}
\end{equation}
其中，$\operatorname{PP}_{\mathrm{lm}}$ 为使用用户数据训练的语言模型的困惑度，$\operatorname{PP}_{\mathrm{public}}$ 为公共模型的困惑度，均对 $S_{\mathrm{uniq}}$ 中的各个序列 $w$ 计算。

\subsection{评估基准}
\refcite{ch06-ref-078}{Mireshghallah 等人（2023）}提出 ConfAIde，用来评估大语言模型能否正确判断应否披露可能敏感的信息。它侧重考察模型能否推断在什么情况下不应分享某项知识，而非分析模型知道什么。基准定义了四类复杂程度依次增加的情境。最简单的情境要求模型给一项信息的敏感程度评分，不考虑由谁访问或如何使用。其他三层逐渐变复杂，最后一层要求模型为一段长对话做记录。多人在不同时点加入和退出，讨论不同种类的信息，最后由模型决定哪些信息可以在面向全组的摘要中分享。

研究者为每个样例收集人工标注，确定信息的敏感程度，以及人在特定语境中对分享或不分享的偏好。

ConfAIde 的作者用这套基准评估了多个表现先进的大语言模型，包括 GPT-4 和 ChatGPT。模型在最基础的敏感信息识别任务上表现尚可，但需要更多理解上下文时，能力明显下降。结果表明，大语言模型难以区分私人信息与公开信息。一般而言，较少分享敏感数据的模型，也更可能隐藏非敏感数据。ConfAIde 表明，即使经过指令调优的大语言模型，仍难以就隐私和保密作推理，可能需要新的技术来弥补这个不足。

\subsection{缓解策略}
本节讨论减轻大语言模型隐私问题的实际办法，分为训练时可采用的方法与推理时可采用的方法。

\subsubsection{训练时的隐私保护}
防止大语言模型散播个人信息，最直观的办法可能是从训练数据中清除这些信息。模型没有接触过私人数据，就不会记住它。如第~\ref{ch:pretraining}~章所述，这是大语言模型预训练中广泛采用的预处理步骤。不过，由于涉及的数据量很大，用常规匿名化技术几乎无法保证已经移除所有 PII。

前面讨论的差分隐私也可用作缓解措施。\refcite{ch06-ref-001}{Abadi 等人（2016）}等研究者将差分隐私纳入随机梯度下降优化器，从而引入训练过程。这个办法有其优点，但迄今在多数情况下会影响训练，通常使模型质量下降。

训练数据匿名化与差分隐私还有一个局限。即使模型没有直接学到个人信息，也可能推断出来。\refcite{ch06-ref-104}{Staab 等人（2023）}发现，几个表现先进的大语言模型可以依据 Reddit 用户发帖的内容，较准确地发现其个人信息。该研究用直接提示的方法推断年龄、性别、位置等个人属性。研究者向各模型提供用户的帖子，要求猜测这些信息。即使将输入数据匿名化，移除用户明确披露个人信息的内容，模型也常能猜对。在作者整理的评估数据集上，GPT-4 的准确率最高，达到 84.6\%。

猜测一位身份不明的 Reddit 用户的大致年龄，乍看可能无害，但很多参加网络论坛的人认为，只要不透露姓名，他们就是匿名的。人们在互联网上也更愿意说出原本不想分享的事。然而，如果他们透露较多关于住处、工作、家庭和年龄的信息，有意追查的人就可能通过社交媒体和公开记录拼凑出其身份。没有大语言模型时也存在这种风险，只是操作较费力。大语言模型可能加快这种恶意活动，使之更容易扩大规模。

未来或可开发更有效的技术，防止模型记住训练数据中的 PII。不过，怎样在持续开发更强模型的同时，阻止它学到足以推断口语或书面语中地域和代际差异的知识，则更难想象。无论其利弊如何，大语言模型已经拥有这种能力。与限制模型能力相比，重点制定使用模型时的防护措施，可能更为审慎。

\subsubsection{推理时的隐私保护}
发现利用模型的新手法后，往往会有很大压力要求处理相应的查询方式。当有人报告直接查询大语言模型可以获得本不应提供的信息后，用户常会很快发现相近提示已经失效。例如，系统可以在内部更新，在上下文中加入指令，要求不要执行确定某人位置的请求。随后，模型可以回答无法提供答案。模型提供者愿意并能够尽快解决这类问题固然有益，但这是事后应对，难以彻底消除隐私问题。

现代大语言模型似乎难以完全保证私人数据得到保护，因此应用开发者也必须考虑这些模型可能怎样使用户面临风险。原书举例说，当时一些大语言模型提供者，如 OpenAI，会保存通过 API 发送的查询，以支持后续技术改进。使用大语言模型的人往往需要考虑，在将提示发送到第三方服务之前，先将其匿名化，尤其当用户可能有意或无意地在应用中输入个人数据时。OpaquePrompts 等工具可以自动移除敏感信息。视使用场景而定，如果有必要，还可以在模型给出回答后，将匿名化时替换掉的内容还原到输出中。

\begin{practicaltip}
对依赖外部开发模型的组织，另一种常见选择是使用开源大语言模型，替代 ChatGPT 一类服务。可以在组织内部部署一份模型，让用户的提示保留在组织网络内，避免将其暴露给第三方服务。这大幅降低了泄露敏感数据的可能，也使部署复杂得多。
\end{practicaltip}

大语言模型需要昂贵的计算资源，要控制这些计算成本，需要专门的人力和知识。除了在第三方大语言模型上构建的应用越来越多，人们也有很强的需求，想直接用 ChatGPT 等产品帮助处理日常事务。不过，许多人担心自己的私人数据暴露，因而有所犹豫。想用在线大语言模型把电子邮件改得更专业的用户，必须把邮件内容提供给模型服务商。为避免这种情况，可以在个人设备上运行的较小模型得到更广泛的使用。

\begin{practicaltip}
原书介绍，Nomic AI 的 GPT4All 是一种常用的可本地安装的大语言模型界面。它提供 Falcon 和 Mistral 等不同模型变体供下载，模型的参数量低于 100 亿。这些模型足够小，可以在个人设备上快速给出质量较好的回答，无需发出 API 请求。与能力更强的 GPT 模型相比，它们也有局限，尤其在需要处理较长上下文时。不过，对回答问题或协助完成基本任务，较小的大语言模型可能已绰绰有余。在许多情况下，为显著减少隐私风险而接受这一折中是合理的。
\end{practicaltip}

硬件行业也密切关注模型本地运行的趋势。过去十年，大多数硬件进展都着眼于在越来越大的数据集上更高效地训练模型。不过，近几年面向神经模型优化的边缘计算吸引了大量投资。有预测认为，这项技术的增长潜力甚至可能超过 NVIDIA 依靠大型 GPU 已取得的收入增长。除了其他因素，大语言模型的隐私问题也推动了人们对去中心化模型的兴趣。

% PDF 290--293；原书第 262--265 页。
\section{教程——衡量与缓解大语言模型的偏见}
\label{sec:ch06-tutorial}
\subsection{概述}
第~\ref{sec:ch06-bias}~节讨论了大语言模型偏见的影响，以及为缓解偏见开发的一些技术。在这个教程中，我们将应用其中一种方法，观察模型行为的相应变化。这个练习参考了 \refcite{ch06-ref-075}{Meade 等人（2022）}的工作。他们考察了多种偏见缓解技术，并在 GitHub 仓库中提供了全部实验的代码。

\begin{keypoint}
\noindent\textbf{目标}
\begin{itemize}
\item 分析 CrowS 基准是怎样设计来衡量偏见的。
\item 在 RoBERTa 上试用一种偏见缓解技术，评估改善程度。
\item 在一项下游任务中应用去偏模型，考察其语言模型能力是否下降。
\end{itemize}
\end{keypoint}

这里给出教程的简明版本。完整版见 \url{https://github.com/springer-llms-deep-dive/llms-deep-dive-tutorials}。

\subsection{实验设计}
本练习将演示如何用 \texttt{bias-bench} 库减少 RoBERTa 模型表现出的性别偏见。随后，用 CrowS 指标展示改善程度，并在情感分析任务中比较去偏模型与原始模型的能力。

CrowS 基准的数据集由成对的句子组成。在每一对中，一个句子体现刻板印象，另一个句子替换相关词语，使之与刻板印象相反。例如，如果是种族刻板印象，可将“black”替换为“white”。如果是性别刻板印象，则可将“woman”替换为“man”，依此类推。除了这些词元，两句话的其他部分完全相同。用这些数据可衡量给定大语言模型的偏见，以及各种偏见缓解技术的相对效果。

本实验选用 Sent-Debias 算法。它的出发点是，如果模型对性别等属性完全中立，那么“He was a slow runner”和“She was a slow runner”的嵌入一般很接近，甚至完全相同。嵌入之间的差异主要可归因于偏见。Sent-Debias 在许多样例中提取这种差异，通过主成分分析将它们映射到较低维的子空间，得到一组表示偏见方向的向量。学得子空间后，将它纳入前向传播。在返回最终输出之前，从任意文本表示中减去它在偏见子空间上的投影。

Sent-Debias 需要一个规模较大、内容多样的数据集，来生成上述流程所用的句子。它有一组预定义的词用于增强数据，例如“boy”和“girl”。我们取维基百科的一批文本来学习模型偏见的表示，通过替换可能带有偏见的词元，从所得句子的嵌入差异中提取偏见信息。

采用偏见缓解方法并评估性别偏见是否比原模型减少后，再比较模型在下游任务中接受微调的能力。这里使用 SST，它是 GLUE 基准中的标准情感分析数据集（\refcite{ch06-ref-102}{Socher et al., 2013}；\refcite{ch06-ref-113}{Wang et al., 2019}）。

\subsection{结果与分析}
开始去偏实验前，先评估 \texttt{roberta-base} 模型在 CrowS 上的表现。该指标反映在可能有偏见的句子中填入被遮盖词元时，模型选择刻板印象的可能性。在模型倾向刻板选项的情况下，数值越低越好。%
\footnote{译注：按本节后文的比较，50.0 表示模型对两类选项无偏好。分数低于 50.0 时，不能再按越低越接近无偏好来解读。}
能力较弱的语言模型也可能在该指标上接近理想表现，因为它对数据中偏见的学习不足，未能选出体现刻板印象的词元。按该指标衡量，能力较弱的大语言模型经常显得比能力更强的模型偏见更少。

\noindent\begin{minipage}{\linewidth}
\begin{lstlisting}[caption={RoBERTa 的 CrowS 指标},label={lst:ch06-crows-original}]
# Output
-----------------------------------------------------------
Total examples: 261
Metric score: 60.15
\end{lstlisting}
\end{minipage}\par

接着用 Sent-Debias 为 RoBERTa 计算性别偏见子空间。得到偏见方向后，再用 CrowS 基准评估性别偏见是否减少。结果已更接近 50.0，说明模型偏好刻板词元的频率低于原来。

\noindent\begin{minipage}{\linewidth}
\begin{lstlisting}[caption={去偏后重新评估 CrowS 指标},label={lst:ch06-crows-debiased}]
# Output
-----------------------------------------------------------
Total examples: 261
Metric score: 52.11
\end{lstlisting}
\end{minipage}\par

\begin{table}[!htbp]
\centering
\caption{比较不同模型在 CrowS 和 SST 基准上的表现，展示去偏的影响。}
\label{tab:ch06-crows-sst}
\begin{tabular}{@{}lrr@{}}
\toprule
模型 & CrowS & SST\\
\midrule
基础 RoBERTa & 60.15 & 0.922\\
采用 Sent-Debias 的 RoBERTa & 52.11 & 0.930\\
\bottomrule
\end{tabular}
\end{table}

如上所述，模型在 CrowS 指标上的改善，可能与整体词元预测能力的下降有关。为确认减轻性别偏见后，大语言模型的任务表现仍然相近，对采用和未采用 Sent-Debias 的模型分别做情感分析微调，比较结果。重复运行训练过程时，数值有小幅变化，但无论是否采用 Sent-Debias，SST 测试数据上的准确率大致相同，见表~\ref{tab:ch06-crows-sst}。

这些结果是积极的，但尚无法确定去偏是否已经成功，或经过处理的模型是否在微调过程中重新学到了一定程度的性别偏见。情感分析训练数据可能本身带有偏见，所用的 CrowS 指标未必容易反映它的影响。还需要更细致地分析这项任务，才能确定缓解偏见的尝试是否成功。

\subsection{结论}
本教程展示了一种有前景的缓解大语言模型偏见的方法，现有技术仍无法彻底解决这一问题。\refcite{ch06-ref-075}{Meade 等人（2022）}的一个重要发现是，虽然已有多种去偏策略，却没有一种能在不同模型和不同类型的偏见上始终有良好表现。他们还发现，用同一算法重复运行时，CrowS、StereoSet 和 SEAT 等基准的评估结果可能有较大波动。这使人们需要进一步研究，这些指标是否足以形成全面的偏见评估。衡量与缓解偏见都还需要进一步研究。

% PDF 293--302。
\begin{chapterreferences}
\urldef\RefUrlB\url{https://en.wikipedia.org/wiki/Wikipedia:Who_writes_Wikipedia%3F}
\refentry{ch06-ref-001}{Martin Abadi, Andy Chu, Ian Goodfellow, H. Brendan McMahan, Ilya Mironov, Kunal Talwar, and Li Zhang. Deep learning with differential privacy. In \textit{Proceedings of the 2016 ACM SIGSAC Conference on Computer and Communications} \textit{Security}, CCS’16. ACM, October 2016. doi: 10.1145/2976749.2978318. URL \url{http://dx.doi.org/10.1145/2976749.2978318}.}
\refentry{ch06-ref-002}{Abubakar Abid, Maheen Farooqi, and James Zou. Persistent anti-muslim bias in large language models. In \textit{Proceedings of the 2021 AAAI/ACM Conference on AI,} \textit{Ethics, and Society}, pages 298–306, 2021.}
\refentry{ch06-ref-003}{Jaimeen Ahn and Alice Oh. Mitigating language-dependent ethnic bias in bert. \textit{arXiv} \textit{preprint arXiv:2109.05704}, 2021.}
\refentry{ch06-ref-004}{Giuseppe Attanasio, Debora Nozza, Dirk Hovy, and Elena Baralis. Entropy-based attention regularization frees unintended bias mitigation from lists. In Smaranda Muresan, Preslav Nakov, and Aline Villavicencio, editors, \textit{Findings of the Association for Computational Linguistics: ACL 2022}, pages 1105–1119, Dublin, Ireland, May 2022. Association for Computational Linguistics. doi: 10.18653/v1/2022.findings-acl.88. URL \url{https://aclanthology.org/2022.findings-acl.88}.}
\refentry{ch06-ref-005}{Emily M Bender and Alexander Koller. Climbing towards nlu: On meaning, form, and understanding in the age of data. In \textit{Proceedings of the 58th annual meeting} \textit{of the association for computational linguistics}, pages 5185–5198, 2020.}
\refentry{ch06-ref-006}{Emily M Bender, Timnit Gebru, Angelina McMillan-Major, and Shmargaret Shmitchell. On the dangers of stochastic parrots: Can language models be too big? In \textit{Proceedings of the 2021 ACM conference on fairness, accountability, and} \textit{transparency}, pages 610–623, 2021.}
\refentry{ch06-ref-007}{Camiel J Beukeboom and Christian Burgers. How stereotypes are shared through language: a review and introduction of the aocial categories and stereotypes communication (scsc) framework. \textit{Review of Communication Research}, 7:1–37, 2019.}
\refentry{ch06-ref-008}{Su Lin Blodgett and Brendan O’Connor. Racial disparity in natural language processing: A case study of social media african-american english. \textit{arXiv preprint} \textit{arXiv:1707.00061}, 2017.}
\refentry{ch06-ref-009}{Rishi Bommasani, Percy Liang, and Tony Lee. Holistic evaluation of language models. \textit{Annals of the New York Academy of Sciences}, 2023.}
\refentry{ch06-ref-010}{Conrad Borchers, Dalia Sara Gala, Benjamin Gilburt, Eduard Oravkin, Wilfried Bounsi, Yuki M Asano, and Hannah Rose Kirk. Looking for a handsome carpenter! debiasing gpt-3 job advertisements. \textit{arXiv preprint arXiv:2205.11374}, 2022.}
\refentry{ch06-ref-011}{Shikha Bordia and Samuel R Bowman. Identifying and reducing gender bias in word-level language models. \textit{arXiv preprint arXiv:1904.03035}, 2019.}
\refentry{ch06-ref-012}{Tom Brown, Benjamin Mann, Nick Ryder, Melanie Subbiah, Jared D Kaplan, Prafulla Dhariwal, Arvind Neelakantan, Pranav Shyam, Girish Sastry, Amanda Askell, et al. Language models are few-shot learners. \textit{Advances in neural information processing systems}, 33:1877–1901, 2020.}
\refentry{ch06-ref-013}{Deng Cai, Yan Wang, Huayang Li, Wai Lam, and Lemao Liu. Neural machine translation with monolingual translation memory. \textit{arXiv preprint arXiv:2105.11269}, 2021.}
\refentry{ch06-ref-014}{Meng Cao, Yue Dong, Jiapeng Wu, and Jackie Chi Kit Cheung. Factual error correction for abstractive summarization models. \textit{arXiv preprint arXiv:2010.08712}, 2020.}
\refentry{ch06-ref-015}{Nicholas Carlini et al. Extracting training data from large language models, 2021.}
\refentry{ch06-ref-016}{Kai-Wei Chang, Vinodkumar Prabhakaran, and Vicente Ordonez. Bias and fairness in natural language processing. In \textit{Proceedings of the 2019 Conference on Empirical Methods in Natural Language Processing and the 9th International Joint Conference on Natural Language Processing (EMNLP-IJCNLP): Tutorial Abstracts}, 2019.}
\refentry{ch06-ref-017}{I Chern, Steffi Chern, Shiqi Chen, Weizhe Yuan, Kehua Feng, Chunting Zhou, Junxian He, Graham Neubig, Pengfei Liu, et al. Factool: Factuality detection in generative ai–a tool augmented framework for multi-task and multi-domain scenarios. \textit{arXiv preprint arXiv:2307.13528}, 2023.}
\refentry{ch06-ref-018}{Aakanksha Chowdhery et al. Palm: Scaling language modeling with pathways, 2022.}
\refentry{ch06-ref-019}{Sumanth Dathathri, Andrea Madotto, Janice Lan, Jane Hung, Eric Frank, Piero Molino, Jason Yosinski, and Rosanne Liu. Plug and play language models: A simple approach to controlled text generation. \textit{arXiv preprint arXiv:1912.02164}, 2019.}
\refentry{ch06-ref-020}{Ameet Deshpande, Vishvak Murahari, Tanmay Rajpurohit, Ashwin Kalyan, and Karthik Narasimhan. Toxicity in chatgpt: Analyzing persona-assigned language models. \textit{arXiv preprint arXiv:2304.05335}, 2023.}
\refentry{ch06-ref-021}{Harnoor Dhingra, Preetiha Jayashanker, Sayali Moghe, and Emma Strubell. Queer people are people first: Deconstructing sexual identity stereotypes in large language models. \textit{arXiv preprint arXiv:2307.00101}, 2023.}
\refentry{ch06-ref-022}{Emily Dinan, Angela Fan, Adina Williams, Jack Urbanek, Douwe Kiela, and Jason Weston. Queens are powerful too: Mitigating gender bias in dialogue generation. \textit{arXiv preprint arXiv:1911.03842}, 2019.}
\refentry{ch06-ref-023}{Lucas Dixon, John Li, Jeffrey Sorensen, Nithum Thain, and Lucy Vasserman. Measuring and mitigating unintended bias in text classification. In \textit{Proceedings of the} \textit{2018 AAAI/ACM Conference on AI, Ethics, and Society}, pages 67–73, 2018.}
\refentry{ch06-ref-024}{Nouha Dziri, Sivan Milton, Mo Yu, Osmar Zaiane, and Siva Reddy. On the origin of hallucinations in conversational models: Is it the datasets or the models? \textit{arXiv} \textit{preprint arXiv:2204.07931}, 2022.}
\refentry{ch06-ref-025}{Zahra Fatemi, Chen Xing, Wenhao Liu, and Caiming Xiong. Improving gender fairness of pre-trained language models without catastrophic forgetting. \textit{arXiv} \textit{preprint arXiv:2110.05367}, 2021.}
\refentry{ch06-ref-026}{Chao Feng, Xinyu Zhang, and Zichu Fei. Knowledge solver: Teaching llms to search for domain knowledge from knowledge graphs. \textit{arXiv preprint arXiv:2309.03118}, 2023.}
\refentry{ch06-ref-027}{Emilio Ferrara. Should chatgpt be biased? challenges and risks of bias in large language models. \textit{arXiv preprint arXiv:2304.03738}, 2023.}
\refentry{ch06-ref-028}{Jessica Ficler and Yoav Goldberg. Controlling linguistic style aspects in neural language generation. \textit{arXiv preprint arXiv:1707.02633}, 2017.}
\refentry{ch06-ref-029}{Isabel O Gallegos, Ryan A Rossi, Joe Barrow, Md Mehrab Tanjim, Sungchul Kim, Franck Dernoncourt, Tong Yu, Ruiyi Zhang, and Nesreen K Ahmed. Bias and fairness in large language models: A survey. \textit{arXiv preprint arXiv:2309.00770}, 2023.}
\refentry{ch06-ref-030}{Leo Gao et al. The pile: An 800gb dataset of diverse text for language modeling, 2020.}
\refentry{ch06-ref-031}{Claire Gardent, Anastasia Shimorina, Shashi Narayan, and Laura Perez-Beltrachini. Creating training corpora for nlg micro-planning. In \textit{55th annual meeting of the} \textit{Association for Computational Linguistics (ACL)}, 2017.}
\refentry{ch06-ref-032}{Aparna Garimella, Rada Mihalcea, and Akhash Amarnath. Demographic-aware language model fine-tuning as a bias mitigation technique. In \textit{Proceedings of the 2nd} \textit{Conference of the Asia-Pacific Chapter of the Association for Computational Linguistics and the 12th International Joint Conference on Natural Language Processing}, pages 311–319, 2022.}
\refentry{ch06-ref-033}{Samuel Gehman, Suchin Gururangan, Maarten Sap, Yejin Choi, and Noah A Smith. Realtoxicityprompts: Evaluating neural toxic degeneration in language models. \textit{arXiv preprint arXiv:2009.11462}, 2020.}
\refentry{ch06-ref-034}{Sayan Ghosh, Mathieu Chollet, Eugene Laksana, Louis-Philippe Morency, and Stefan Scherer. Affect-lm: A neural language model for customizable affective text generation. \textit{arXiv preprint arXiv:1704.06851}, 2017.}
\refentry{ch06-ref-035}{Michael Gira, Ruisu Zhang, and Kangwook Lee. Debiasing pre-trained language models via efficient fine-tuning. In \textit{Proceedings of the Second Workshop on Language Technology for Equality, Diversity and Inclusion}, pages 59–69, 2022.}
\refentry{ch06-ref-036}{Seraphina Goldfarb-Tarrant, Rebecca Marchant, Ricardo Muñoz Sánchez, Mugdha Pandya, and Adam Lopez. Intrinsic bias metrics do not correlate with application bias. \textit{arXiv preprint arXiv:2012.15859}, 2020.}
\refentry{ch06-ref-037}{Zhibin Gou, Zhihong Shao, Yeyun Gong, Yelong Shen, Yujiu Yang, Nan Duan, and Weizhu Chen. Critic: Large language models can self-correct with tool-interactive critiquing. \textit{arXiv preprint arXiv:2305.11738}, 2023.}
\refentry{ch06-ref-038}{Anthony G Greenwald, Debbie E McGhee, and Jordan LK Schwartz. Measuring individual differences in implicit cognition: the implicit association test. \textit{Journal} \textit{of personality and social psychology}, 74(6):1464, 1998.}
\refentry{ch06-ref-039}{Chuan Guo, Geoff Pleiss, Yu Sun, and Kilian Q Weinberger. On calibration of modern neural networks. In \textit{International conference on machine learning}, pages 1321–1330. PMLR, 2017.}
\refentry{ch06-ref-040}{Demi Guo, Alexander M Rush, and Yoon Kim. Parameter-efficient transfer learning with diff pruning. \textit{arXiv preprint arXiv:2012.07463}, 2020.}
\refentry{ch06-ref-041}{Wei Guo and Aylin Caliskan. Detecting emergent intersectional biases: Contextualized word embeddings contain a distribution of human-like biases. In \textit{Proceedings} \textit{of the 2021 AAAI/ACM Conference on AI, Ethics, and Society}, pages 122–133, 2021.}
\refentry{ch06-ref-042}{Suchin Gururangan, Ana Marasović, Swabha Swayamdipta, Kyle Lo, Iz Beltagy, Doug Downey, and Noah A. Smith. Don’t stop pretraining: Adapt language models to domains and tasks, 2020.}
\refentry{ch06-ref-043}{Xudong Han, Timothy Baldwin, and Trevor Cohn. Towards equal opportunity fairness through adversarial learning. \textit{arXiv preprint arXiv:2203.06317}, 2022.}
\refentry{ch06-ref-044}{Lukas Hauzenberger, Shahed Masoudian, Deepak Kumar, Markus Schedl, and Navid Rekabsaz. Modular and on-demand bias mitigation with attribute-removal subnetworks. In \textit{Findings of the Association for Computational Linguistics: ACL} \textit{2023}, pages 6192–6214, 2023.}
\refentry{ch06-ref-045}{Daniel Hershcovich, Stella Frank, Heather Lent, Miryam de Lhoneux, Mostafa Abdou, Stephanie Brandl, Emanuele Bugliarello, Laura Cabello Piqueras, Ilias Chalkidis, Ruixiang Cui, et al. Challenges and strategies in cross-cultural nlp. \textit{arXiv preprint arXiv:2203.10020}, 2022.}
\refentry{ch06-ref-046}{Ari Holtzman, Jan Buys, Li Du, Maxwell Forbes, and Yejin Choi. The curious case of neural text degeneration. \textit{arXiv preprint arXiv:1904.09751}, 2019.}
\refentry{ch06-ref-047}{Jeremy Howard and Sebastian Ruder. Universal language model fine-tuning for text classification, 2018.}
\refentry{ch06-ref-048}{Huseyin A. Inan, Osman Ramadan, Lukas Wutschitz, Daniel Jones, Victor Rühle, James Withers, and Robert Sim. Training data leakage analysis in language models, 2021.}
\refentry{ch06-ref-049}{Nishtha Jain, Maja Popovic, Declan Groves, and Eva Vanmassenhove. Generating gender augmented data for nlp. \textit{arXiv preprint arXiv:2107.05987}, 2021.}
\refentry{ch06-ref-050}{Przemyslaw Joniak and Akiko Aizawa. Gender biases and where to find them: Exploring gender bias in pre-trained transformer-based language models using movement pruning. \textit{arXiv preprint arXiv:2207.02463}, 2022.}
\refentry{ch06-ref-051}{Saurav Kadavath, Tom Conerly, Amanda Askell, Tom Henighan, Dawn Drain, Ethan Perez, Nicholas Schiefer, Zac Hatfield-Dodds, Nova DasSarma, Eli Tran-Johnson, et al. Language models (mostly) know what they know. \textit{arXiv preprint} \textit{arXiv:2207.05221}, 2022.}
\refentry{ch06-ref-052}{Nitish Shirish Keskar, Bryan McCann, Lav R Varshney, Caiming Xiong, and Richard Socher. Ctrl: A conditional transformer language model for controllable generation. \textit{arXiv preprint arXiv:1909.05858}, 2019.}
\refentry{ch06-ref-053}{Siwon Kim, Sangdoo Yun, Hwaran Lee, Martin Gubri, Sungroh Yoon, and Seong Joon Oh. Propile: Probing privacy leakage in large language models, 2023.}
\refentry{ch06-ref-054}{Ben Krause, Akhilesh Deepak Gotmare, Bryan McCann, Nitish Shirish Keskar, Shafiq Joty, Richard Socher, and Nazneen Fatema Rajani. Gedi: Generative discriminator guided sequence generation. \textit{arXiv preprint arXiv:2009.06367}, 2020.}
\refentry{ch06-ref-055}{Julia Kreutzer, Isaac Caswell, Lisa Wang, Ahsan Wahab, Daan van Esch, Nasanbayar Ulzii-Orshikh, Allahsera Tapo, Nishant Subramani, Artem Sokolov, Claytone Sikasote, et al. Quality at a glance: An audit of web-crawled multilingual datasets. \textit{Transactions of the Association for Computational Linguistics}, 10:50–72, 2022.}
\refentry{ch06-ref-056}{Keita Kurita, Anna Belova, and Antonios Anastasopoulos. Towards robust toxic content classification. \textit{arXiv preprint arXiv:1912.06872}, 2019.}
\refentry{ch06-ref-057}{Anne Lauscher, Tobias Lueken, and Goran Glavaš. Sustainable modular debiasing of language models. \textit{arXiv preprint arXiv:2109.03646}, 2021.}
\refentry{ch06-ref-058}{Nayeon Lee, Wei Ping, Peng Xu, Mostofa Patwary, Pascale N Fung, Mohammad Shoeybi, and Bryan Catanzaro. Factuality enhanced language models for open-ended text generation. \textit{Advances in Neural Information Processing Systems}, 35: 34586–34599, 2022.}
\refentry{ch06-ref-059}{Alyssa Lees, Vinh Q Tran, Yi Tay, Jeffrey Sorensen, Jai Gupta, Donald Metzler, and Lucy Vasserman. A new generation of perspective api: Efficient multilingual character-level transformers. In \textit{Proceedings of the 28th ACM SIGKDD Conference on Knowledge Discovery and Data Mining}, pages 3197–3207, 2022.}
\refentry{ch06-ref-060}{Junyi Li, Xiaoxue Cheng, Wayne Xin Zhao, Jian-Yun Nie, and Ji-Rong Wen. Halueval: A large-scale hallucination evaluation benchmark for large language models. In \textit{Proceedings of the 2023 Conference on Empirical Methods in Natural Language Processing}, pages 6449–6464, 2023a.}
\refentry{ch06-ref-061}{Shaobo Li, Xiaoguang Li, Lifeng Shang, Zhenhua Dong, Chengjie Sun, Bingquan Liu, Zhenzhou Ji, Xin Jiang, and Qun Liu. How pre-trained language models capture factual knowledge? a causal-inspired analysis. \textit{arXiv preprint} \textit{arXiv:2203.16747}, 2022.}
\refentry{ch06-ref-062}{Xingxuan Li, Ruochen Zhao, Yew Ken Chia, Bosheng Ding, Lidong Bing, Shafiq Joty, and Soujanya Poria. Chain of knowledge: A framework for grounding large language models with structured knowledge bases. \textit{arXiv preprint} \textit{arXiv:2305.13269}, 2023b.}
\refentry{ch06-ref-063}{Yuanzhi Li, Sébastien Bubeck, Ronen Eldan, Allie Del Giorno, Suriya Gunasekar, and Yin Tat Lee. Textbooks are all you need ii: phi-1.5 technical report. \textit{arXiv} \textit{preprint arXiv:2309.05463}, 2023c.}
\refentry{ch06-ref-064}{Stephanie Lin, Jacob Hilton, and Owain Evans. Truthfulqa: Measuring how models mimic human falsehoods. \textit{arXiv preprint arXiv:2109.07958}, 2021.}
\refentry{ch06-ref-065}{Adam Liska, Tomas Kocisky, Elena Gribovskaya, Tayfun Terzi, Eren Sezener, Devang Agrawal, D’Autume Cyprien De Masson, Tim Scholtes, Manzil Zaheer, Susannah Young, et al. Streamingqa: A benchmark for adaptation to new knowledge over time in question answering models. In \textit{International Conference on Machine} \textit{Learning}, pages 13604–13622. PMLR, 2022.}
\refentry{ch06-ref-066}{Haochen Liu, Da Tang, Ji Yang, Xiangyu Zhao, Hui Liu, Jiliang Tang, and Youlong Cheng. Rating distribution calibration for selection bias mitigation in recommendations. In \textit{Proceedings of the ACM Web Conference 2022}, WWW ’22, page 2048–2057, New York, NY, USA, 2022. Association for Computing Machinery. ISBN 9781450390965. doi: 10.1145/3485447.3512078. URL \url{https://doi.org/10.1145/3485447.3512078}.}
\refentry{ch06-ref-067}{Pengfei Liu, Weizhe Yuan, Jinlan Fu, Zhengbao Jiang, Hiroaki Hayashi, and Graham Neubig. Pre-train, prompt, and predict: A systematic survey of prompting methods in natural language processing. \textit{ACM Computing Surveys}, 55(9):1–35, 2023.}
\refentry{ch06-ref-068}{Kaiji Lu, Piotr Mardziel, Fangjing Wu, Preetam Amancharla, and Anupam Datta. Gender bias in neural natural language processing. \textit{Logic, Language, and Security:} \textit{Essays Dedicated to Andre Scedrov on the Occasion of His 65th Birthday}, pages 189–202, 2020.}
\refentry{ch06-ref-069}{Ziyang Luo, Can Xu, Pu Zhao, Xiubo Geng, Chongyang Tao, Jing Ma, Qingwei Lin, and Daxin Jiang. Augmented large language models with parametric knowledge guiding. \textit{arXiv preprint arXiv:2305.04757}, 2023.}
\refentry{ch06-ref-070}{Potsawee Manakul, Adian Liusie, and Mark JF Gales. Selfcheckgpt: Zero-resource black-box hallucination detection for generative large language models. \textit{arXiv} \textit{preprint arXiv:2303.08896}, 2023.}
\refentry{ch06-ref-071}{Justus Mattern, Zhijing Jin, Mrinmaya Sachan, Rada Mihalcea, and Bernhard Schölkopf. Understanding stereotypes in language models: Towards robust measurement and zero-shot debiasing. \textit{arXiv preprint arXiv:2212.10678}, 2022.}
\refentry{ch06-ref-072}{Rowan Hall Maudslay, Hila Gonen, Ryan Cotterell, and Simone Teufel. It’s all in the name: Mitigating gender bias with name-based counterfactual data substitution. \textit{arXiv preprint arXiv:1909.00871}, 2019.}
\refentry{ch06-ref-073}{Chandler May, Alex Wang, Shikha Bordia, Samuel R Bowman, and Rachel Rudinger. On measuring social biases in sentence encoders. \textit{arXiv preprint} \textit{arXiv:1903.10561}, 2019.}
\refentry{ch06-ref-074}{Nick McKenna, Tianyi Li, Liang Cheng, Mohammad Javad Hosseini, Mark Johnson, and Mark Steedman. Sources of hallucination by large language models on inference tasks. \textit{arXiv preprint arXiv:2305.14552}, 2023.}
\refentry{ch06-ref-075}{Nicholas Meade, Elinor Poole-Dayan, and Siva Reddy. An empirical survey of the effectiveness of debiasing techniques for pre-trained language models. In \textit{Proceedings of the 60th Annual Meeting of the Association for Computational Linguistics (Volume 1: Long Papers)}, pages 1878–1898, Dublin, Ireland, May 2022. Association for Computational Linguistics. doi: 10.18653/v1/2022.acl-long.132. URL \url{https://aclanthology.org/2022.acl-long.132}.}
\refentry{ch06-ref-076}{Grégoire Mialon, Roberto Dessì, Maria Lomeli, Christoforos Nalmpantis, Ram Pasunuru, Roberta Raileanu, Baptiste Rozière, Timo Schick, Jane Dwivedi-Yu, Asli Celikyilmaz, Edouard Grave, Yann LeCun, and Thomas Scialom. Augmented language models: a survey, 2023.}
\refentry{ch06-ref-077}{Sewon Min, Kalpesh Krishna, Xinxi Lyu, Mike Lewis, Wen-tau Yih, Pang Wei Koh, Mohit Iyyer, Luke Zettlemoyer, and Hannaneh Hajishirzi. Factscore: Fine-grained atomic evaluation of factual precision in long form text generation. \textit{arXiv preprint} \textit{arXiv:2305.14251}, 2023.}
\refentry{ch06-ref-078}{Niloofar Mireshghallah, Hyunwoo Kim, Xuhui Zhou, Yulia Tsvetkov, Maarten Sap, Reza Shokri, and Yejin Choi. Can llms keep a secret? testing privacy implications of language models via contextual integrity theory. \textit{arXiv preprint} \textit{arXiv:2310.17884}, 2023.}
\refentry{ch06-ref-079}{Niels Mündler, Jingxuan He, Slobodan Jenko, and Martin Vechev. Self-contradictory hallucinations of large language models: Evaluation, detection and mitigation. \textit{arXiv preprint arXiv:2305.15852}, 2023.}
\refentry{ch06-ref-080}{Moin Nadeem, Anna Bethke, and Siva Reddy. Stereoset: Measuring stereotypical bias in pretrained language models. \textit{arXiv preprint arXiv:2004.09456}, 2020.}
\refentry{ch06-ref-081}{Nikita Nangia, Clara Vania, Rasika Bhalerao, and Samuel R Bowman. Crows-pairs: A challenge dataset for measuring social biases in masked language models. \textit{arXiv} \textit{preprint arXiv:2010.00133}, 2020.}
\refentry{ch06-ref-082}{Milad Nasr, Nicholas Carlini, Jonathan Hayase, Matthew Jagielski, A. Feder Cooper, Daphne Ippolito, Christopher A. Choquette-Choo, Eric Wallace, Florian Tramèr, and Katherine Lee. Scalable extraction of training data from (production) language models, 2023.}
\refentry{ch06-ref-083}{Roberto Navigli and Simone Paolo Ponzetto. Babelnet: The automatic construction, evaluation and application of a wide-coverage multilingual semantic network. \textit{Artificial intelligence}, 193:217–250, 2012.}
\refentry{ch06-ref-084}{Roberto Navigli, Simone Conia, and Björn Ross. Biases in large language models: Origins, inventory and discussion. \textit{ACM Journal of Data and Information Quality}, 2023.}
\refentry{ch06-ref-085}{Debora Nozza, Federico Bianchi, Dirk Hovy, et al. Honest: Measuring hurtful sentence completion in language models. In \textit{Proceedings of the 2021 Conference} \textit{of the North American Chapter of the Association for Computational Linguistics:} \textit{Human Language Technologies}. Association for Computational Linguistics, 2021.}
\refentry{ch06-ref-086}{Ali Omrani, Alireza Salkhordeh Ziabari, Charles Yu, Preni Golazizian, Brendan Kennedy, Mohammad Atari, Heng Ji, and Morteza Dehghani. Social-group-agnostic bias mitigation via the stereotype content model. In \textit{Proc. The 61st Annual Meeting of the Association for Computational Linguistics (ACL2023)}, 2023.}
\refentry{ch06-ref-087}{Hadas Orgad and Yonatan Belinkov. Blind: Bias removal with no demographics. In \textit{Proceedings of the 61st Annual Meeting of the Association for Computational} \textit{Linguistics (Volume 1: Long Papers)}, pages 8801–8821, 2023.}
\refentry{ch06-ref-088}{SunYoung Park, Kyuri Choi, Haeun Yu, and Youngjoong Ko. Never too late to learn: Regularizing gender bias in coreference resolution. In \textit{Proceedings of the} \textit{Sixteenth ACM International Conference on Web Search and Data Mining}, pages 15–23, 2023.}
\refentry{ch06-ref-089}{John Pavlopoulos, Jeffrey Sorensen, Lucas Dixon, Nithum Thain, and Ion Androutsopoulos. Toxicity detection: Does context really matter? \textit{arXiv preprint} \textit{arXiv:2006.00998}, 2020.}
\refentry{ch06-ref-090}{Guilherme Penedo, Quentin Malartic, Daniel Hesslow, Ruxandra Cojocaru, Alessandro Cappelli, Hamza Alobeidli, Baptiste Pannier, Ebtesam Almazrouei, and Julien Launay. The refinedweb dataset for falcon llm: outperforming curated corpora with web data, and web data only. \textit{arXiv preprint arXiv:2306.01116}, 2023.}
\refentry{ch06-ref-091}{Ethan Perez, Sam Ringer, Kamilė Lukošiūtė, Karina Nguyen, Edwin Chen, Scott Heiner, Craig Pettit, Catherine Olsson, Sandipan Kundu, Saurav Kadavath, et al. Discovering language model behaviors with model-written evaluations. \textit{arXiv} \textit{preprint arXiv:2212.09251}, 2022.}
\refentry{ch06-ref-092}{Yusu Qian, Urwa Muaz, Ben Zhang, and Jae Won Hyun. Reducing gender bias in word-level language models with a gender-equalizing loss function. \textit{arXiv preprint} \textit{arXiv:1905.12801}, 2019.}
\refentry{ch06-ref-093}{Pranav Rajpurkar, Jian Zhang, Konstantin Lopyrev, and Percy Liang. Squad: 100,000+ questions for machine comprehension of text, 2016.}
\refentry{ch06-ref-094}{Leonardo Ranaldi, Elena Sofia Ruzzetti, Davide Venditti, Dario Onorati, and Fabio Massimo Zanzotto. A trip towards fairness: Bias and de-biasing in large language models. \textit{arXiv preprint arXiv:2305.13862}, 2023.}
\refentry{ch06-ref-095}{Ruiyang Ren, Yuhao Wang, Yingqi Qu, Wayne Xin Zhao, Jing Liu, Hao Tian, Hua Wu, Ji-Rong Wen, and Haifeng Wang. Investigating the factual knowledge boundary of large language models with retrieval augmentation, 2023.}
\refentry{ch06-ref-096}{Stephen Robertson, Hugo Zaragoza, et al. The probabilistic relevance framework: Bm25 and beyond. \textit{Foundations and Trends® in Information Retrieval}, 3(4):333–389, 2009.}
\refentry{ch06-ref-097}{Yash Savani, Colin White, and Naveen Sundar Govindarajulu. Intra-processing methods for debiasing neural networks, 2020.}
\refentry{ch06-ref-098}{Timo Schick, Sahana Udupa, and Hinrich Schütze. Self-diagnosis and self-debiasing: A proposal for reducing corpus-based bias in nlp, 2021.}
\refentry{ch06-ref-099}{John Schulman. Reinforcement learning from human feedback: progress and challenges. In \textit{Berkley Electrical Engineering and Computer Sciences. URL:} \url{https://eecs.berkeley.edu/research/colloquium/230419} [accessed 2023-11-15], 2023.}
\refentry{ch06-ref-100}{Weijia Shi, Sewon Min, Michihiro Yasunaga, Minjoon Seo, Rich James, Mike Lewis, Luke Zettlemoyer, and Wen-tau Yih. Replug: Retrieval-augmented black-box language models. \textit{arXiv preprint arXiv:2301.12652}, 2023.}
\refentry{ch06-ref-101}{Eric Michael Smith, Melissa Hall, Melanie Kambadur, Eleonora Presani, and Adina Williams. “i’m sorry to hear that”: Finding new biases in language models with a holistic descriptor dataset. In \textit{Proceedings of the 2022 Conference on Empirical} \textit{Methods in Natural Language Processing}, pages 9180–9211, 2022.}
\refentry{ch06-ref-102}{Richard Socher, Alex Perelygin, Jean Wu, Jason Chuang, Christopher D. Manning, Andrew Ng, and Christopher Potts. Recursive deep models for semantic compositionality over a sentiment treebank. In \textit{Proceedings of the 2013 Conference} \textit{on Empirical Methods in Natural Language Processing}, pages 1631–1642, Seattle, Washington, USA, October 2013. Association for Computational Linguistics. URL \url{https://www.aclweb.org/anthology/D13-1170}.}
\refentry{ch06-ref-103}{Irene Solaiman and Christy Dennison. Process for adapting language models to society (palms) with values-targeted datasets. \textit{Advances in Neural Information} \textit{Processing Systems}, 34:5861–5873, 2021.}
\refentry{ch06-ref-104}{Robin Staab, Mark Vero, Mislav Balunović, and Martin Vechev. Beyond memorization: Violating privacy via inference with large language models, 2023.}
\refentry{ch06-ref-105}{Tianxiang Sun, Xiaotian Zhang, Zhengfu He, Peng Li, Qinyuan Cheng, Hang Yan, Xiangyang Liu, Yunfan Shao, Qiong Tang, Xingjian Zhao, et al. Moss: Training conversational language models from synthetic data. \textit{arXiv preprint} \textit{arXiv:2307.15020}, 7, 2023.}
\refentry{ch06-ref-106}{Tony Sun, Kellie Webster, Apu Shah, William Yang Wang, and Melvin Johnson. They, them, theirs: Rewriting with gender-neutral english, 2021.}
\refentry{ch06-ref-107}{Himanshu Thakur, Atishay Jain, Praneetha Vaddamanu, Paul Pu Liang, and Louis-Philippe Morency. Language models get a gender makeover: Mitigating gender bias with few-shot data interventions. \textit{arXiv preprint arXiv:2306.04597}, 2023.}
\refentry{ch06-ref-108}{Ewoenam Kwaku Tokpo and Toon Calders. Text style transfer for bias mitigation using masked language modeling. \textit{arXiv preprint arXiv:2201.08643}, 2022.}
\refentry{ch06-ref-109}{Hugo Touvron et al. Llama 2: Open foundation and fine-tuned chat models, 2023.}
\refentry{ch06-ref-110}{Prasetya Ajie Utama, Nafise Sadat Moosavi, and Iryna Gurevych. Towards debiasing nlu models from unknown biases. \textit{arXiv preprint arXiv:2009.12303}, 2020.}
\refentry{ch06-ref-111}{Pranav Narayanan Venkit, Sanjana Gautam, Ruchi Panchanadikar, Shomir Wilson, et al. Nationality bias in text generation. \textit{arXiv preprint arXiv:2302.02463}, 2023.}
\refentry{ch06-ref-112}{Eric Wallace, Tony Z Zhao, Shi Feng, and Sameer Singh. Concealed data poisoning attacks on nlp models. \textit{arXiv preprint arXiv:2010.12563}, 2020.}
\refentry{ch06-ref-113}{Alex Wang, Amanpreet Singh, Julian Michael, Felix Hill, Omer Levy, and Samuel R. Bowman. Glue: A multi-task benchmark and analysis platform for natural language understanding, 2019.}
\refentry{ch06-ref-114}{Xinyi Wang, Sebastian Ruder, and Graham Neubig. Expanding pretrained models to thousands more languages via lexicon-based adaptation. \textit{arXiv preprint} \textit{arXiv:2203.09435}, 2022.}
\refentry{ch06-ref-115}{Kellie Webster, Xuezhi Wang, Ian Tenney, Alex Beutel, Emily Pitler, Ellie Pavlick, Jilin Chen, Ed Chi, and Slav Petrov. Measuring and reducing gendered correlations in pre-trained models. \textit{arXiv preprint arXiv:2010.06032}, 2020.}
\refentry{ch06-ref-116}{Jason Wei, Xuezhi Wang, Dale Schuurmans, Maarten Bosma, Fei Xia, Ed Chi, Quoc V Le, Denny Zhou, et al. Chain-of-thought prompting elicits reasoning in large language models. \textit{Advances in Neural Information Processing Systems}, 35:24824–24837, 2022.}
\refentry{ch06-ref-117}{Wikipedia Contributors. Who writes wikipedia? \RefUrlB{}, 2023. Accessed: 2023-04-01.}
\refentry{ch06-ref-118}{Weiqi Wu, Chengyue Jiang, Yong Jiang, Pengjun Xie, and Kewei Tu. Do plms know and understand ontological knowledge? \textit{arXiv preprint arXiv:2309.05936}, 2023.}
\refentry{ch06-ref-119}{Miao Xiong, Zhiyuan Hu, Xinyang Lu, Yifei Li, Jie Fu, Junxian He, and Bryan Hooi. Can llms express their uncertainty? an empirical evaluation of confidence elicitation in llms. \textit{arXiv preprint arXiv:2306.13063}, 2023.}
\refentry{ch06-ref-120}{Ke Yang, Charles Yu, Yi R Fung, Manling Li, and Heng Ji. Adept: A debiasing prompt framework. In \textit{Proceedings of the AAAI Conference on Artificial Intelligence}, volume 37, pages 10780–10788, 2023.}
\refentry{ch06-ref-121}{Shunyu Yao, Jeffrey Zhao, Dian Yu, Nan Du, Izhak Shafran, Karthik Narasimhan, and Yuan Cao. React: Synergizing reasoning and acting in language models. \textit{arXiv} \textit{preprint arXiv:2210.03629}, 2022.}
\refentry{ch06-ref-122}{Zhangyue Yin, Qiushi Sun, Qipeng Guo, Jiawen Wu, Xipeng Qiu, and Xuanjing Huang. Do large language models know what they don’t know? \textit{arXiv preprint} \textit{arXiv:2305.18153}, 2023.}
\refentry{ch06-ref-123}{Charles Yu, Sullam Jeoung, Anish Kasi, Pengfei Yu, and Heng Ji. Unlearning bias in language models by partitioning gradients. In \textit{Findings of the Association for} \textit{Computational Linguistics: ACL 2023}, pages 6032–6048, 2023.}
\refentry{ch06-ref-124}{Sina Zarrieß, Henrik Voigt, and Simeon Schüz. Decoding methods in neural language generation: a survey. \textit{Information}, 12(9):355, 2021.}
\refentry{ch06-ref-125}{Abdelrahman Zayed, Goncalo Mordido, Samira Shabanian, and Sarath Chandar. Should we attend more or less? modulating attention for fairness. \textit{arXiv preprint} \textit{arXiv:2305.13088}, 2023.}
\refentry{ch06-ref-126}{Yuheng Zha, Yichi Yang, Ruichen Li, and Zhiting Hu. Alignscore: Evaluating factual consistency with a unified alignment function. \textit{arXiv preprint arXiv:2305.16739}, 2023.}
\refentry{ch06-ref-127}{Hongyi Zhang, Moustapha Cisse, Yann N Dauphin, and David Lopez-Paz. mixup: Beyond empirical risk minimization. \textit{arXiv preprint arXiv:1710.09412}, 2017.}
\refentry{ch06-ref-128}{Susan Zhang, Stephen Roller, Naman Goyal, Mikel Artetxe, Moya Chen, Shuohui Chen, Christopher Dewan, Mona Diab, Xian Li, Xi Victoria Lin, et al. Opt: Open pre-trained transformer language models. \textit{arXiv preprint arXiv:2205.01068}, 2022.}
\refentry{ch06-ref-129}{Yue Zhang, Yafu Li, Leyang Cui, Deng Cai, Lemao Liu, Tingchen Fu, Xinting Huang, Enbo Zhao, Yu Zhang, Yulong Chen, et al. Siren’s song in the ai ocean: A survey on hallucination in large language models. \textit{arXiv preprint} \textit{arXiv:2309.01219}, 2023.}
\refentry{ch06-ref-130}{Ruochen Zhao, Xingxuan Li, Shafiq Joty, Chengwei Qin, and Lidong Bing. Verify-and-edit: A knowledge-enhanced chain-of-thought framework. \textit{arXiv preprint} \textit{arXiv:2305.03268}, 2023.}
\refentry{ch06-ref-131}{Wayne Xin Zhao, Jing Liu, Ruiyang Ren, and Ji-Rong Wen. Dense text retrieval based on pretrained language models: A survey. \textit{arXiv preprint arXiv:2211.14876}, 2022.}
\refentry{ch06-ref-132}{Chunting Zhou, Pengfei Liu, Puxin Xu, Srini Iyer, Jiao Sun, Yuning Mao, Xuezhe Ma, Avia Efrat, Ping Yu, Lili Yu, et al. Lima: Less is more for alignment. \textit{arXiv} \textit{preprint arXiv:2305.11206}, 2023.}
\end{chapterreferences}

